\documentclass[../main.tex]{subfiles}
\begin{document}

\chapter{I/O Management}
\lessoninfo{
  video={P3L5},
  textbook={OSTEP \cite{ostep} ch.~36--37, 39--42; Silberschatz et al.\ \cite{silberschatz2018} ch.~12; Kerrisk \cite{kerrisk2010} ch.~13--14},
  keywords={device register, device driver, block device, character device, memory-mapped I/O, polling, programmed I/O, DMA, operating system bypass, virtual file system, inode, ext2, buffer cache, journaling},
}

Besides the CPU and memory, whose management Lessons~7 and~8 described, an
operating system manages the devices through which a computer receives
input, produces output and stores data persistently.  This lesson examines
how the operating system represents and controls such devices: the
interface a device presents, the drivers that encapsulate device-specific
knowledge, and the ways in which the CPU and a device exchange commands,
data and notifications.  It then follows a file access down the block
device stack and studies how file systems---the virtual file system and ext2
in particular---organize data on disk and reduce the cost of accessing it.

\section{I/O devices}
\label{sec:l11-devices}

Input/output (I/O) devices connect a computer to its environment.\source{P3L5,
lesson preview, I/O devices, basic I/O device features; OSTEP ch.~36;
Silberschatz et al.\ ch.~12.}  They differ widely in function---some only
deliver input to the computer, some only accept output from it, many do
both---and in speed, from a device that produces a few characters per second
to a network interface that transfers gigabits per second.  The I/O
management of an operating system follows three principles:
\begin{itemize}
  \item it defines \emph{protocols}, standard interfaces through which
    devices are accessed;
  \item it provides \emph{dedicated handlers}, the device drivers and
    interrupt handlers, for the device-specific work;
  \item it \emph{decouples} the details of individual devices from the rest
    of the system, so that applications use abstractions such as files
    rather than device operations.
\end{itemize}

\begin{definition}[Device register]
  A \term{device register}[デバイスレジスタ] is a register of a device that
  the CPU reads or writes to interact with the device.  Devices typically
  provide \emph{command} registers, which the CPU writes to tell the device
  what to do; \emph{data} registers, through which data are transferred to
  and from the device; and \emph{status} registers, which the CPU reads to
  learn what the device is doing.
\end{definition}

The registers form the interface of the device (\cref{fig:l11-device}).
Behind it lie the device internals: a microcontroller, which is the device's
own CPU, on-device memory, and other device-specific chips, such as the
analog-to-digital converters of a microphone.  Only the interface is
visible to the operating system; how the device uses its internals to carry
out a command is hidden from it.

\begin{figure}[htbp]
  \centering
  % A basic I/O device: an interface of device registers that the CPU reads
% and writes, and device-specific internals hidden behind it.
% Usage: % A basic I/O device: an interface of device registers that the CPU reads
% and writes, and device-specific internals hidden behind it.
% Usage: % A basic I/O device: an interface of device registers that the CPU reads
% and writes, and device-specific internals hidden behind it.
% Usage: \input{figures/l11-device}   (width ~116 mm)
\begin{tikzpicture}[x=1mm, y=1mm, font=\footnotesize\sffamily,
  >={Stealth[length=1.6mm]},
  reg/.style={draw, fill=white, minimum width=24mm, minimum height=6mm,
              inner sep=1pt},
  part/.style={draw, fill=white, minimum width=44mm, minimum height=6mm,
               inner sep=1pt, font=\scriptsize\sffamily},
  lab/.style={font=\scriptsize\sffamily, align=center},
  hdr/.style={font=\scriptsize\sffamily\bfseries, text=ln-accent}]
  % device outline
  \draw[fill=ln-box] (30,-13) rectangle (118,19.5);
  \node[hdr] at (50,15.5) {\iflangja{インタフェース（デバイスレジスタ）}{interface (device registers)}};
  \node[hdr] at (94,15.5) {\iflangja{内部構成}{internals}};
  \draw[densely dashed, ln-muted] (69,-10.5) -- (69,12);
  \node[reg] (st) at (50,7) {\iflangja{ステータス}{status}};
  \node[reg] (cm) at (50,0) {\iflangja{コマンド}{command}};
  \node[reg] (dt) at (50,-7) {\iflangja{データ}{data}};
  \node[part] (mc) at (94,7) {\iflangja{マイクロコントローラ（デバイスの CPU）}{microcontroller (the device's CPU)}};
  \node[part] (mm) at (94,0) {\iflangja{デバイス上のメモリ}{on-device memory}};
  \node[part] (ch) at (94,-7) {\iflangja{その他の回路（A/D 変換器など）}{other chips (e.g.\ A/D converter)}};
  \draw[<->] (cm.east) -- (mm.west);
  % CPU side
  \node[draw, fill=ln-accent!22, minimum width=12mm, minimum height=8mm]
    (cpu) at (7,0) {CPU};
  \draw[<->, thick] (cpu.east) -- (34,0);
  \node[lab, anchor=south] at (22,0.6) {\iflangja{読み書き}{read /}};
  \node[lab, anchor=north] at (22,-0.6) {\iflangja{}{write}};
  \draw (34,7) -- (34,-7);
  \foreach \y in {7,0,-7} \draw[->] (34,\y) -- (38,\y);
\end{tikzpicture}
   (width ~116 mm)
\begin{tikzpicture}[x=1mm, y=1mm, font=\footnotesize\sffamily,
  >={Stealth[length=1.6mm]},
  reg/.style={draw, fill=white, minimum width=24mm, minimum height=6mm,
              inner sep=1pt},
  part/.style={draw, fill=white, minimum width=44mm, minimum height=6mm,
               inner sep=1pt, font=\scriptsize\sffamily},
  lab/.style={font=\scriptsize\sffamily, align=center},
  hdr/.style={font=\scriptsize\sffamily\bfseries, text=ln-accent}]
  % device outline
  \draw[fill=ln-box] (30,-13) rectangle (118,19.5);
  \node[hdr] at (50,15.5) {\iflangja{インタフェース（デバイスレジスタ）}{interface (device registers)}};
  \node[hdr] at (94,15.5) {\iflangja{内部構成}{internals}};
  \draw[densely dashed, ln-muted] (69,-10.5) -- (69,12);
  \node[reg] (st) at (50,7) {\iflangja{ステータス}{status}};
  \node[reg] (cm) at (50,0) {\iflangja{コマンド}{command}};
  \node[reg] (dt) at (50,-7) {\iflangja{データ}{data}};
  \node[part] (mc) at (94,7) {\iflangja{マイクロコントローラ（デバイスの CPU）}{microcontroller (the device's CPU)}};
  \node[part] (mm) at (94,0) {\iflangja{デバイス上のメモリ}{on-device memory}};
  \node[part] (ch) at (94,-7) {\iflangja{その他の回路（A/D 変換器など）}{other chips (e.g.\ A/D converter)}};
  \draw[<->] (cm.east) -- (mm.west);
  % CPU side
  \node[draw, fill=ln-accent!22, minimum width=12mm, minimum height=8mm]
    (cpu) at (7,0) {CPU};
  \draw[<->, thick] (cpu.east) -- (34,0);
  \node[lab, anchor=south] at (22,0.6) {\iflangja{読み書き}{read /}};
  \node[lab, anchor=north] at (22,-0.6) {\iflangja{}{write}};
  \draw (34,7) -- (34,-7);
  \foreach \y in {7,0,-7} \draw[->] (34,\y) -- (38,\y);
\end{tikzpicture}
   (width ~116 mm)
\begin{tikzpicture}[x=1mm, y=1mm, font=\footnotesize\sffamily,
  >={Stealth[length=1.6mm]},
  reg/.style={draw, fill=white, minimum width=24mm, minimum height=6mm,
              inner sep=1pt},
  part/.style={draw, fill=white, minimum width=44mm, minimum height=6mm,
               inner sep=1pt, font=\scriptsize\sffamily},
  lab/.style={font=\scriptsize\sffamily, align=center},
  hdr/.style={font=\scriptsize\sffamily\bfseries, text=ln-accent}]
  % device outline
  \draw[fill=ln-box] (30,-13) rectangle (118,19.5);
  \node[hdr] at (50,15.5) {\iflangja{インタフェース（デバイスレジスタ）}{interface (device registers)}};
  \node[hdr] at (94,15.5) {\iflangja{内部構成}{internals}};
  \draw[densely dashed, ln-muted] (69,-10.5) -- (69,12);
  \node[reg] (st) at (50,7) {\iflangja{ステータス}{status}};
  \node[reg] (cm) at (50,0) {\iflangja{コマンド}{command}};
  \node[reg] (dt) at (50,-7) {\iflangja{データ}{data}};
  \node[part] (mc) at (94,7) {\iflangja{マイクロコントローラ（デバイスの CPU）}{microcontroller (the device's CPU)}};
  \node[part] (mm) at (94,0) {\iflangja{デバイス上のメモリ}{on-device memory}};
  \node[part] (ch) at (94,-7) {\iflangja{その他の回路（A/D 変換器など）}{other chips (e.g.\ A/D converter)}};
  \draw[<->] (cm.east) -- (mm.west);
  % CPU side
  \node[draw, fill=ln-accent!22, minimum width=12mm, minimum height=8mm]
    (cpu) at (7,0) {CPU};
  \draw[<->, thick] (cpu.east) -- (34,0);
  \node[lab, anchor=south] at (22,0.6) {\iflangja{読み書き}{read /}};
  \node[lab, anchor=north] at (22,-0.6) {\iflangja{}{write}};
  \draw (34,7) -- (34,-7);
  \foreach \y in {7,0,-7} \draw[->] (34,\y) -- (38,\y);
\end{tikzpicture}

  \caption{A basic I/O device.  The CPU interacts with the device only
    through its command, status and data registers; the internals that carry
    out the commands are device-specific.}
  \label{fig:l11-device}
\end{figure}

\section{CPU--device interconnect}
\label{sec:l11-interconnect}

Devices are attached to the CPU and memory through controllers on an
interconnect.\source{P3L5, CPU device interconnect; OSTEP ch.~36;
Silberschatz et al.\ ch.~12.}  The interconnect limits how fast data can
move between a device and the rest of the system and how many devices can
be attached.  A standard family of such interconnects is the \term{Peripheral Component Interconnect}[PCI]%
\index{PCI|see{Peripheral Component Interconnect}} (PCI).  Its original form
is a shared parallel bus, which PCI-X (PCI eXtended) made faster.  Current
platforms use PCI Express (PCIe), which remains compatible with the software
interface of PCI but offers more bandwidth, faster and lower-latency access,
and connections for more devices.\margin{Although commonly called a bus,
PCIe connects each device through a point-to-point serial link and
switches.}

Not every device is connected to PCIe directly (\cref{fig:l11-interconnect}).
A SCSI controller attached to PCIe connects disks on its own SCSI bus, and
an expansion bus interface connects slower peripherals such as keyboards on
an expansion (peripheral) bus.  Bridging controllers handle the differences
between the buses, and the bridge that connects the CPU to memory also
connects both to PCIe.  An access from the CPU to a device may therefore
cross several interconnects.

\begin{figure}[htbp]
  \centering
  % Typical interconnect: CPU and memory behind a bridge/memory controller,
% device controllers on PCIe, and further buses behind their controllers.
% Usage: % Typical interconnect: CPU and memory behind a bridge/memory controller,
% device controllers on PCIe, and further buses behind their controllers.
% Usage: % Typical interconnect: CPU and memory behind a bridge/memory controller,
% device controllers on PCIe, and further buses behind their controllers.
% Usage: \input{figures/l11-interconnect}   (width ~116 mm)
\begin{tikzpicture}[x=1mm, y=1mm, font=\scriptsize\sffamily,
  box/.style={draw, fill=white, minimum height=6mm, inner sep=1pt, align=center},
  ctl/.style={draw, fill=ln-box, minimum height=8.5mm, minimum width=25mm,
              inner sep=1pt, align=center},
  bus/.style={line width=1.6pt},
  blab/.style={font=\scriptsize\sffamily\itshape, text=ln-muted}]
  % top row
  \node[box, fill=ln-accent!22, minimum width=18mm, minimum height=8.5mm]
    (cpu) at (16,27) {CPU};
  \node[ctl, minimum width=34mm] (br) at (58,27)
    {\iflangja{ブリッジ／\\メモリコントローラ}{bridge /\\memory controller}};
  \node[box, minimum width=22mm, minimum height=8.5mm] (mem) at (100,27)
    {\iflangja{メモリ}{memory}};
  \draw[thick] (cpu) -- (br);
  \draw[thick] (br) -- (mem);
  % PCIe
  \draw[bus] (3,14) -- (113,14);
  \node[blab, anchor=south west] at (3,14.6) {PCIe};
  \draw[thick] (br.south) -- (58,14);
  % controllers
  \node[ctl] (gc) at (16,3) {\iflangja{グラフィックス\\コントローラ}{graphics\\controller}};
  \node[ctl] (ni) at (44,3) {\iflangja{ネットワーク\\インタフェース}{network\\interface}};
  \node[ctl] (sc) at (72,3) {\iflangja{SCSI\\コントローラ}{SCSI\\controller}};
  \node[ctl] (ex) at (100,3) {\iflangja{拡張バス\\インタフェース}{expansion bus\\interface}};
  \foreach \n in {gc,ni,sc,ex} \draw[thick] (\n.north) -- (\n.north |- 0,14);
  % below the controllers
  \node[box, minimum width=18mm] (disp) at (16,-15) {\iflangja{ディスプレイ}{display}};
  \draw (gc.south) -- (disp.north);
  \node[box, minimum width=18mm] (net) at (44,-15) {\iflangja{ネットワーク}{network}};
  \draw (ni.south) -- (net.north);
  \draw[bus] (59.5,-9) -- (84.5,-9);
  \node[blab, anchor=south west] at (59.5,-8.6) {\iflangja{SCSI バス}{SCSI bus}};
  \draw (80,-1.25) -- (80,-9);
  \node[box, minimum width=10mm] (d1) at (66,-17) {\iflangja{ディスク}{disk}};
  \node[box, minimum width=10mm] (d2) at (78,-17) {\iflangja{ディスク}{disk}};
  \draw (d1.north) -- (66,-9);  \draw (d2.north) -- (78,-9);
  \draw[bus] (87.5,-9) -- (112.5,-9);
  \node[blab, anchor=south west] at (87.5,-8.6) {\iflangja{拡張バス}{expansion bus}};
  \draw (109,-1.25) -- (109,-9);
  \node[box, minimum width=13mm] (kb) at (94,-17) {\iflangja{キーボード}{keyboard}};
  \node[box, minimum width=10mm] (ms) at (107,-17) {\iflangja{マウス}{mouse}};
  \draw (kb.north) -- (94,-9);  \draw (ms.north) -- (107,-9);
\end{tikzpicture}
   (width ~116 mm)
\begin{tikzpicture}[x=1mm, y=1mm, font=\scriptsize\sffamily,
  box/.style={draw, fill=white, minimum height=6mm, inner sep=1pt, align=center},
  ctl/.style={draw, fill=ln-box, minimum height=8.5mm, minimum width=25mm,
              inner sep=1pt, align=center},
  bus/.style={line width=1.6pt},
  blab/.style={font=\scriptsize\sffamily\itshape, text=ln-muted}]
  % top row
  \node[box, fill=ln-accent!22, minimum width=18mm, minimum height=8.5mm]
    (cpu) at (16,27) {CPU};
  \node[ctl, minimum width=34mm] (br) at (58,27)
    {\iflangja{ブリッジ／\\メモリコントローラ}{bridge /\\memory controller}};
  \node[box, minimum width=22mm, minimum height=8.5mm] (mem) at (100,27)
    {\iflangja{メモリ}{memory}};
  \draw[thick] (cpu) -- (br);
  \draw[thick] (br) -- (mem);
  % PCIe
  \draw[bus] (3,14) -- (113,14);
  \node[blab, anchor=south west] at (3,14.6) {PCIe};
  \draw[thick] (br.south) -- (58,14);
  % controllers
  \node[ctl] (gc) at (16,3) {\iflangja{グラフィックス\\コントローラ}{graphics\\controller}};
  \node[ctl] (ni) at (44,3) {\iflangja{ネットワーク\\インタフェース}{network\\interface}};
  \node[ctl] (sc) at (72,3) {\iflangja{SCSI\\コントローラ}{SCSI\\controller}};
  \node[ctl] (ex) at (100,3) {\iflangja{拡張バス\\インタフェース}{expansion bus\\interface}};
  \foreach \n in {gc,ni,sc,ex} \draw[thick] (\n.north) -- (\n.north |- 0,14);
  % below the controllers
  \node[box, minimum width=18mm] (disp) at (16,-15) {\iflangja{ディスプレイ}{display}};
  \draw (gc.south) -- (disp.north);
  \node[box, minimum width=18mm] (net) at (44,-15) {\iflangja{ネットワーク}{network}};
  \draw (ni.south) -- (net.north);
  \draw[bus] (59.5,-9) -- (84.5,-9);
  \node[blab, anchor=south west] at (59.5,-8.6) {\iflangja{SCSI バス}{SCSI bus}};
  \draw (80,-1.25) -- (80,-9);
  \node[box, minimum width=10mm] (d1) at (66,-17) {\iflangja{ディスク}{disk}};
  \node[box, minimum width=10mm] (d2) at (78,-17) {\iflangja{ディスク}{disk}};
  \draw (d1.north) -- (66,-9);  \draw (d2.north) -- (78,-9);
  \draw[bus] (87.5,-9) -- (112.5,-9);
  \node[blab, anchor=south west] at (87.5,-8.6) {\iflangja{拡張バス}{expansion bus}};
  \draw (109,-1.25) -- (109,-9);
  \node[box, minimum width=13mm] (kb) at (94,-17) {\iflangja{キーボード}{keyboard}};
  \node[box, minimum width=10mm] (ms) at (107,-17) {\iflangja{マウス}{mouse}};
  \draw (kb.north) -- (94,-9);  \draw (ms.north) -- (107,-9);
\end{tikzpicture}
   (width ~116 mm)
\begin{tikzpicture}[x=1mm, y=1mm, font=\scriptsize\sffamily,
  box/.style={draw, fill=white, minimum height=6mm, inner sep=1pt, align=center},
  ctl/.style={draw, fill=ln-box, minimum height=8.5mm, minimum width=25mm,
              inner sep=1pt, align=center},
  bus/.style={line width=1.6pt},
  blab/.style={font=\scriptsize\sffamily\itshape, text=ln-muted}]
  % top row
  \node[box, fill=ln-accent!22, minimum width=18mm, minimum height=8.5mm]
    (cpu) at (16,27) {CPU};
  \node[ctl, minimum width=34mm] (br) at (58,27)
    {\iflangja{ブリッジ／\\メモリコントローラ}{bridge /\\memory controller}};
  \node[box, minimum width=22mm, minimum height=8.5mm] (mem) at (100,27)
    {\iflangja{メモリ}{memory}};
  \draw[thick] (cpu) -- (br);
  \draw[thick] (br) -- (mem);
  % PCIe
  \draw[bus] (3,14) -- (113,14);
  \node[blab, anchor=south west] at (3,14.6) {PCIe};
  \draw[thick] (br.south) -- (58,14);
  % controllers
  \node[ctl] (gc) at (16,3) {\iflangja{グラフィックス\\コントローラ}{graphics\\controller}};
  \node[ctl] (ni) at (44,3) {\iflangja{ネットワーク\\インタフェース}{network\\interface}};
  \node[ctl] (sc) at (72,3) {\iflangja{SCSI\\コントローラ}{SCSI\\controller}};
  \node[ctl] (ex) at (100,3) {\iflangja{拡張バス\\インタフェース}{expansion bus\\interface}};
  \foreach \n in {gc,ni,sc,ex} \draw[thick] (\n.north) -- (\n.north |- 0,14);
  % below the controllers
  \node[box, minimum width=18mm] (disp) at (16,-15) {\iflangja{ディスプレイ}{display}};
  \draw (gc.south) -- (disp.north);
  \node[box, minimum width=18mm] (net) at (44,-15) {\iflangja{ネットワーク}{network}};
  \draw (ni.south) -- (net.north);
  \draw[bus] (59.5,-9) -- (84.5,-9);
  \node[blab, anchor=south west] at (59.5,-8.6) {\iflangja{SCSI バス}{SCSI bus}};
  \draw (80,-1.25) -- (80,-9);
  \node[box, minimum width=10mm] (d1) at (66,-17) {\iflangja{ディスク}{disk}};
  \node[box, minimum width=10mm] (d2) at (78,-17) {\iflangja{ディスク}{disk}};
  \draw (d1.north) -- (66,-9);  \draw (d2.north) -- (78,-9);
  \draw[bus] (87.5,-9) -- (112.5,-9);
  \node[blab, anchor=south west] at (87.5,-8.6) {\iflangja{拡張バス}{expansion bus}};
  \draw (109,-1.25) -- (109,-9);
  \node[box, minimum width=13mm] (kb) at (94,-17) {\iflangja{キーボード}{keyboard}};
  \node[box, minimum width=10mm] (ms) at (107,-17) {\iflangja{マウス}{mouse}};
  \draw (kb.north) -- (94,-9);  \draw (ms.north) -- (107,-9);
\end{tikzpicture}

  \caption{A typical platform interconnect.  Device controllers are attached
    to PCIe; disks and keyboards are attached to further buses behind their
    own controllers.  PCIe is drawn as a bus for simplicity; each controller
    is in fact attached through its own point-to-point link.}
  \label{fig:l11-interconnect}
\end{figure}

\section{Device drivers}
\label{sec:l11-drivers}

The operating system does not itself contain the knowledge of how to operate
each device.\source{P3L5, device drivers; OSTEP ch.~36; Silberschatz et
al.\ ch.~12.}

\begin{definition}[Device driver]
  A \term{device driver}[デバイスドライバ] is a device-specific software
  component of the operating system that is responsible for all access to,
  management of and control of the devices of one type.
\end{definition}

There is one driver per device type, and drivers are typically provided by
the manufacturers of the devices; in a modular operating system they are
loaded as kernel modules (Lesson~1).  The operating system in turn
standardizes the interfaces that drivers implement, the \emph{device driver
framework}.  This standardization yields two properties:
\begin{description}
  \item[Device independence] the operating system does not have to be
    specialized for particular devices, since it uses every device of a type
    through the same interface.
  \item[Device diversity] the operating system can support arbitrarily
    different devices, including new ones: a device can be used as soon as
    a driver implements the interface for it.
\end{description}

\section{Types of devices}
\label{sec:l11-types}

Operating systems group devices into classes.\source{P3L5, types of
devices; Silberschatz et al.\ ch.~12; Kerrisk ch.~14.}  The members of a
class are used in the same way, and each class has its own driver
interface.  Three classes are distinguished:
\begin{itemize}
  \item A \term{block device}[ブロックデバイス], such as a disk, stores data
    in fixed-size blocks, and every individual block can be read or written
    directly by its number.
  \item A \term{character device}[キャラクタデバイス], such as a keyboard,
    delivers or accepts a sequential stream of characters; its basic
    operations get or put one character.
  \item A \term{network device}[ネットワークデバイス], such as an Ethernet
    interface, lies between the two: it delivers a stream of data, but in
    chunks of varying size rather than in fixed-size blocks or single
    characters.
\end{itemize}

\subsection{Devices as files}

Operating systems represent devices in the file system.

\begin{definition}[Device file]
  A \term{device file}[デバイスファイル], also called a \emph{special file},
  is a file that represents a device.  The operations that a process
  applies to it---open, read, write, close---are carried out on the device
  by the device's driver.
\end{definition}

In UNIX-like systems device files reside in the directory \texttt{/dev},
which is managed by a special file system such as devfs or, on current
Linux, a tmpfs-based one (tmpfs was the home of POSIX shared memory in
Lesson~9).  A long listing with \texttt{ls -l} marks a block device file
with the type \texttt{b} and a character device file with
\texttt{c}.\margin{On Linux, network devices are the exception: they have no
device files and are used through sockets.}  Because a device looks like a
file, programs that know nothing about devices can use it.

\begin{example}
  The command
  \begin{lstlisting}[numbers=none]
dd if=/dev/sda of=head.img bs=512 count=1
\end{lstlisting}
  copies the first 512~bytes of the first disk into an ordinary file.  The
  program \texttt{dd} uses only open, read and write; that its input is a
  block device, and the disk protocol that retrieves the block, are handled
  by the operating system and the disk's driver.
\end{example}

Not every device file corresponds to hardware.

\begin{definition}[Pseudo device]
  A \term{pseudo device}[疑似デバイス] is a device, represented by a device
  file like any other, that corresponds to no physical hardware; its driver
  implements the device's functionality entirely in software.
\end{definition}

For example, reading \texttt{/dev/zero} returns an endless stream of zero
bytes, which is convenient for creating a file of a given size filled with
zeros.

\begin{keypoint}
  A device is operated through its command, data and status registers, by a
  device driver that hides the device behind an interface the operating
  system standardizes for the whole class of devices.  The classes are block,
  character and network devices, and the first two are represented as device
  files, so that ordinary file operations reach them.
\end{keypoint}

\section{CPU--device interactions}
\label{sec:l11-interactions}

To use a device, the CPU must be able to address its registers, and the
device must be able to tell the CPU when it needs attention.\source{P3L5,
CPU device interactions; OSTEP ch.~36; Silberschatz et al.\ ch.~12.}  There
are two ways to address the registers.

\begin{definition}[Memory-mapped I/O]
  With \term{memory-mapped I/O}[メモリマップドI/O] the device registers
  appear to the CPU as memory locations at specific physical addresses.  A
  part of the physical address space is dedicated to device interactions,
  and the CPU accesses a register with an ordinary load or store instruction
  on its address, which the platform directs to the device instead of to
  memory.
\end{definition}

For a PCI device, the base address registers (BARs) in the device's
configuration space determine how large its region of the physical address
space is and where the region lies; they are configured during the boot
process.  Memory-mapped I/O must not be confused with a memory-mapped file
(Lesson~9), which maps the contents of a file into the virtual address
space of a process.

The second way uses dedicated CPU instructions, such as \texttt{in} and
\texttt{out} on x86.  An \term{I/O port}[I/Oポート] is an address, separate
from the memory addresses, that designates a device register; the
instruction specifies the target port and the value to be written to the
register, or receives the value read from it.

\subsection{Interrupts and polling}

In the other direction, a device must inform the CPU that it has completed
a request or that data have arrived.  It can raise an interrupt (Lesson~5),
or the operating system can find out by polling.

\begin{definition}[Polling]
  The technique in which the operating system reads the status register of
  a device, at times of its own choosing, to find out whether the device
  requires attention, instead of waiting for the device to raise an
  interrupt, is called \term{polling}[ポーリング].
\end{definition}

\Cref{tab:l11-int-poll} compares the two.  An interrupt reaches the CPU as
soon as the device generates it, but every interrupt incurs the handling
overhead: executing the interrupt handler, setting and resetting interrupt
masks, and indirect effects such as the eviction of the interrupted
thread's data from the CPU caches.  Polling happens when it is convenient
for the operating system, but if it happens rarely, events wait until the
next poll; if it happens often, most polls find nothing and waste CPU time.
Which method is preferable depends on the type of device and on the
objectives, such as how frequent the events are and how quickly they must be
handled.

\begin{table}[htbp]
  \centering
  \caption{Interrupts and polling.}
  \label{tab:l11-int-poll}
  \small
  \begin{tabular}{@{}>{\raggedright\arraybackslash}p{0.2\linewidth}>{\raggedright\arraybackslash}p{0.36\linewidth}>{\raggedright\arraybackslash}p{0.36\linewidth}@{}}
    \toprule
    & Interrupt & Polling \\
    \midrule
    Initiated by & the device & the operating system \\
    Event noticed & as soon as it is generated & at the next poll \\
    Cost & handling overhead per interrupt, cache pollution & one status
      read per poll, including polls that find nothing \\
    Drawback & overhead when events are frequent & delay if polls are rare,
      CPU overhead if they are frequent \\
    \bottomrule
  \end{tabular}
\end{table}

\begin{example}
  Suppose that taking an interrupt costs \SI{4}{\micro\second} of CPU time
  beyond the processing of the event itself, and that one poll, which
  collects all events pending since the previous poll, costs
  \SI{0.5}{\micro\second}.  A device that reports 50 events per second costs
  \SI{200}{\micro\second} of CPU time per second with interrupts; polling
  every \SI{1}{\milli\second} would cost \SI{500}{\micro\second} per second
  and delay each event by up to \SI{1}{\milli\second}.  A network interface
  that receives \num{100000} packets per second, however, costs
  \SI{400}{\milli\second} per second---\SI{40}{\percent} of a CPU---with one
  interrupt per packet, while polling every \SI{1}{\milli\second} still costs
  only \SI{500}{\micro\second} per second.
\end{example}

\section{Device access: programmed I/O and DMA}
\label{sec:l11-pio-dma}

The CPU and a device can exchange data in two ways, which differ in whether
the data pass through the CPU.\source{P3L5, device access PIO, device access
DMA; OSTEP ch.~36.}  Both ways use the device registers of
\cref{sec:l11-devices} to issue commands.

\begin{definition}[Programmed I/O]
  With \term{programmed I/O}[プログラム制御I/O]%
  \index{PIO|see{programmed I/O}} (PIO) the CPU performs every step of the
  interaction through the device registers: it writes the command register
  to request an operation, and it moves the data by writing or reading the
  data register, one register-sized piece at a time.
\end{definition}

PIO requires no hardware support beyond the registers.  To send a network
packet with PIO, the CPU writes a command requesting transmission into the
command register, copies the first part of the packet into the data
register, and repeats the copy until the whole packet has been transferred
(\cref{fig:l11-pio-dma}a).  The number of CPU accesses grows with the size of
the data.

\begin{definition}[Direct memory access]
  With \term{direct memory access}[DMA]\index{DMA|see{direct memory access}}
  (DMA) the CPU still issues commands through the command register, but the
  data do not pass through the CPU.  The CPU places the data in a buffer in
  main memory and passes the address and size of the buffer to the device;
  a DMA controller then copies the data between memory and the device on its
  own.
\end{definition}

To send the packet with DMA, the CPU writes the command, places the packet
in a memory buffer, and configures the DMA controller with the buffer's
address and size (\cref{fig:l11-pio-dma}b).  The number of CPU accesses no
longer depends on the size of the data.

\begin{figure}[htbp]
  \centering
  % Sending a packet with programmed I/O (data copied by the CPU into the data
% register) and with DMA (device copies the data from a memory buffer).
% Usage: % Sending a packet with programmed I/O (data copied by the CPU into the data
% register) and with DMA (device copies the data from a memory buffer).
% Usage: % Sending a packet with programmed I/O (data copied by the CPU into the data
% register) and with DMA (device copies the data from a memory buffer).
% Usage: \input{figures/l11-pio-dma}   (width ~118 mm)
\begin{tikzpicture}[x=1mm, y=1mm, font=\footnotesize\sffamily,
  >={Stealth[length=1.6mm]},
  box/.style={draw, fill=white, minimum height=7mm, inner sep=1pt, align=center},
  reg/.style={draw, fill=white, minimum width=20mm, minimum height=5.5mm,
              inner sep=1pt, font=\scriptsize\sffamily},
  num/.style={circle, draw=black, fill=white, text=black, line width=0.4pt,
              inner sep=0.4pt, font=\scriptsize\sffamily\bfseries},
  lab/.style={font=\scriptsize\sffamily, align=left},
  ttl/.style={font=\footnotesize\sffamily\bfseries, anchor=west}]
  % ---------------- (a) PIO
  \node[ttl] at (0,26) {\iflangja{(a) プログラム制御 I/O}{(a) programmed I/O}};
  \node[box, fill=ln-accent!22, minimum width=14mm] (cpu) at (9,13) {CPU};
  \node[box, minimum width=16mm, minimum height=9mm] (mem) at (9,-4)
    {\iflangja{メモリ\\（パケット）}{memory\\(packet)}};
  \draw[fill=ln-box] (30,-9) rectangle (56,21);
  \node[font=\scriptsize\sffamily\bfseries, text=ln-accent, anchor=north]
    at (43,20.5) {NIC};
  \node[reg] (cm) at (43,13) {\iflangja{コマンド}{command}};
  \node[reg] (dt) at (43,3) {\iflangja{データ}{data}};
  \draw[->] (cpu.east) -- node[num, above=0.3mm] {1} (cm.west);
  \draw[->, line width=1.4pt, ln-accent] (mem.north) -- (cpu.south);
  \draw[->, line width=1.4pt, ln-accent] (cpu.south east) -- node[num, below left=-0.4mm] {2} (dt.west);
  \node[lab, anchor=north west, text width=56mm] at (0,-12)
    {\iflangja{\textbf{1} コマンドレジスタへの書き込み 1 回\\
      \textbf{2} データレジスタへの書き込み（パケット全体を\\
      \phantom{\textbf{2} }レジスタの大きさずつ繰り返す）}%
     {\textbf{1} one write to the command register\\
      \textbf{2} writes to the data register, repeated\\
      \phantom{\textbf{2} }until the whole packet is transferred}};
  % ---------------- (b) DMA
  \begin{scope}[xshift=62mm]
  \node[ttl] at (0,26) {\iflangja{(b) DMA}{(b) direct memory access}};
  \node[box, fill=ln-accent!22, minimum width=14mm] (cpu2) at (9,13) {CPU};
  \node[box, minimum width=16mm, minimum height=9mm] (mem2) at (9,-4)
    {\iflangja{メモリ\\（バッファ）}{memory\\(buffer)}};
  \draw[fill=ln-box] (30,-9) rectangle (56,21);
  \node[font=\scriptsize\sffamily\bfseries, text=ln-accent, anchor=north]
    at (43,20.5) {NIC};
  \node[reg] (cm2) at (43,13) {\iflangja{コマンド}{command}};
  \node[reg, fill=ln-accent!10] (dma) at (43,-4) {\iflangja{DMA コントローラ}{DMA controller}};
  \draw[->] (cpu2.east) -- node[num, above=0.3mm] {1} (cm2.west);
  \draw[->] (cpu2.south east) -- node[num, above right=-0.6mm] {2} (dma.north west);
  \draw[->, line width=1.4pt, ln-accent] (mem2.east |- dma.west) -- node[num, below=0.8mm] {3} (dma.west);
  \node[lab, anchor=north west, text width=56mm] at (0,-12)
    {\iflangja{\textbf{1} コマンドレジスタへの書き込み 1 回\\
      \textbf{2} DMA の設定（バッファのアドレスとサイズ）\\
      \textbf{3} デバイスがバッファからデータを読む}%
     {\textbf{1} one write to the command register\\
      \textbf{2} DMA configuration: buffer address and size\\
      \textbf{3} the device reads the data from the buffer}};
  \end{scope}
\end{tikzpicture}
   (width ~118 mm)
\begin{tikzpicture}[x=1mm, y=1mm, font=\footnotesize\sffamily,
  >={Stealth[length=1.6mm]},
  box/.style={draw, fill=white, minimum height=7mm, inner sep=1pt, align=center},
  reg/.style={draw, fill=white, minimum width=20mm, minimum height=5.5mm,
              inner sep=1pt, font=\scriptsize\sffamily},
  num/.style={circle, draw=black, fill=white, text=black, line width=0.4pt,
              inner sep=0.4pt, font=\scriptsize\sffamily\bfseries},
  lab/.style={font=\scriptsize\sffamily, align=left},
  ttl/.style={font=\footnotesize\sffamily\bfseries, anchor=west}]
  % ---------------- (a) PIO
  \node[ttl] at (0,26) {\iflangja{(a) プログラム制御 I/O}{(a) programmed I/O}};
  \node[box, fill=ln-accent!22, minimum width=14mm] (cpu) at (9,13) {CPU};
  \node[box, minimum width=16mm, minimum height=9mm] (mem) at (9,-4)
    {\iflangja{メモリ\\（パケット）}{memory\\(packet)}};
  \draw[fill=ln-box] (30,-9) rectangle (56,21);
  \node[font=\scriptsize\sffamily\bfseries, text=ln-accent, anchor=north]
    at (43,20.5) {NIC};
  \node[reg] (cm) at (43,13) {\iflangja{コマンド}{command}};
  \node[reg] (dt) at (43,3) {\iflangja{データ}{data}};
  \draw[->] (cpu.east) -- node[num, above=0.3mm] {1} (cm.west);
  \draw[->, line width=1.4pt, ln-accent] (mem.north) -- (cpu.south);
  \draw[->, line width=1.4pt, ln-accent] (cpu.south east) -- node[num, below left=-0.4mm] {2} (dt.west);
  \node[lab, anchor=north west, text width=56mm] at (0,-12)
    {\iflangja{\textbf{1} コマンドレジスタへの書き込み 1 回\\
      \textbf{2} データレジスタへの書き込み（パケット全体を\\
      \phantom{\textbf{2} }レジスタの大きさずつ繰り返す）}%
     {\textbf{1} one write to the command register\\
      \textbf{2} writes to the data register, repeated\\
      \phantom{\textbf{2} }until the whole packet is transferred}};
  % ---------------- (b) DMA
  \begin{scope}[xshift=62mm]
  \node[ttl] at (0,26) {\iflangja{(b) DMA}{(b) direct memory access}};
  \node[box, fill=ln-accent!22, minimum width=14mm] (cpu2) at (9,13) {CPU};
  \node[box, minimum width=16mm, minimum height=9mm] (mem2) at (9,-4)
    {\iflangja{メモリ\\（バッファ）}{memory\\(buffer)}};
  \draw[fill=ln-box] (30,-9) rectangle (56,21);
  \node[font=\scriptsize\sffamily\bfseries, text=ln-accent, anchor=north]
    at (43,20.5) {NIC};
  \node[reg] (cm2) at (43,13) {\iflangja{コマンド}{command}};
  \node[reg, fill=ln-accent!10] (dma) at (43,-4) {\iflangja{DMA コントローラ}{DMA controller}};
  \draw[->] (cpu2.east) -- node[num, above=0.3mm] {1} (cm2.west);
  \draw[->] (cpu2.south east) -- node[num, above right=-0.6mm] {2} (dma.north west);
  \draw[->, line width=1.4pt, ln-accent] (mem2.east |- dma.west) -- node[num, below=0.8mm] {3} (dma.west);
  \node[lab, anchor=north west, text width=56mm] at (0,-12)
    {\iflangja{\textbf{1} コマンドレジスタへの書き込み 1 回\\
      \textbf{2} DMA の設定（バッファのアドレスとサイズ）\\
      \textbf{3} デバイスがバッファからデータを読む}%
     {\textbf{1} one write to the command register\\
      \textbf{2} DMA configuration: buffer address and size\\
      \textbf{3} the device reads the data from the buffer}};
  \end{scope}
\end{tikzpicture}
   (width ~118 mm)
\begin{tikzpicture}[x=1mm, y=1mm, font=\footnotesize\sffamily,
  >={Stealth[length=1.6mm]},
  box/.style={draw, fill=white, minimum height=7mm, inner sep=1pt, align=center},
  reg/.style={draw, fill=white, minimum width=20mm, minimum height=5.5mm,
              inner sep=1pt, font=\scriptsize\sffamily},
  num/.style={circle, draw=black, fill=white, text=black, line width=0.4pt,
              inner sep=0.4pt, font=\scriptsize\sffamily\bfseries},
  lab/.style={font=\scriptsize\sffamily, align=left},
  ttl/.style={font=\footnotesize\sffamily\bfseries, anchor=west}]
  % ---------------- (a) PIO
  \node[ttl] at (0,26) {\iflangja{(a) プログラム制御 I/O}{(a) programmed I/O}};
  \node[box, fill=ln-accent!22, minimum width=14mm] (cpu) at (9,13) {CPU};
  \node[box, minimum width=16mm, minimum height=9mm] (mem) at (9,-4)
    {\iflangja{メモリ\\（パケット）}{memory\\(packet)}};
  \draw[fill=ln-box] (30,-9) rectangle (56,21);
  \node[font=\scriptsize\sffamily\bfseries, text=ln-accent, anchor=north]
    at (43,20.5) {NIC};
  \node[reg] (cm) at (43,13) {\iflangja{コマンド}{command}};
  \node[reg] (dt) at (43,3) {\iflangja{データ}{data}};
  \draw[->] (cpu.east) -- node[num, above=0.3mm] {1} (cm.west);
  \draw[->, line width=1.4pt, ln-accent] (mem.north) -- (cpu.south);
  \draw[->, line width=1.4pt, ln-accent] (cpu.south east) -- node[num, below left=-0.4mm] {2} (dt.west);
  \node[lab, anchor=north west, text width=56mm] at (0,-12)
    {\iflangja{\textbf{1} コマンドレジスタへの書き込み 1 回\\
      \textbf{2} データレジスタへの書き込み（パケット全体を\\
      \phantom{\textbf{2} }レジスタの大きさずつ繰り返す）}%
     {\textbf{1} one write to the command register\\
      \textbf{2} writes to the data register, repeated\\
      \phantom{\textbf{2} }until the whole packet is transferred}};
  % ---------------- (b) DMA
  \begin{scope}[xshift=62mm]
  \node[ttl] at (0,26) {\iflangja{(b) DMA}{(b) direct memory access}};
  \node[box, fill=ln-accent!22, minimum width=14mm] (cpu2) at (9,13) {CPU};
  \node[box, minimum width=16mm, minimum height=9mm] (mem2) at (9,-4)
    {\iflangja{メモリ\\（バッファ）}{memory\\(buffer)}};
  \draw[fill=ln-box] (30,-9) rectangle (56,21);
  \node[font=\scriptsize\sffamily\bfseries, text=ln-accent, anchor=north]
    at (43,20.5) {NIC};
  \node[reg] (cm2) at (43,13) {\iflangja{コマンド}{command}};
  \node[reg, fill=ln-accent!10] (dma) at (43,-4) {\iflangja{DMA コントローラ}{DMA controller}};
  \draw[->] (cpu2.east) -- node[num, above=0.3mm] {1} (cm2.west);
  \draw[->] (cpu2.south east) -- node[num, above right=-0.6mm] {2} (dma.north west);
  \draw[->, line width=1.4pt, ln-accent] (mem2.east |- dma.west) -- node[num, below=0.8mm] {3} (dma.west);
  \node[lab, anchor=north west, text width=56mm] at (0,-12)
    {\iflangja{\textbf{1} コマンドレジスタへの書き込み 1 回\\
      \textbf{2} DMA の設定（バッファのアドレスとサイズ）\\
      \textbf{3} デバイスがバッファからデータを読む}%
     {\textbf{1} one write to the command register\\
      \textbf{2} DMA configuration: buffer address and size\\
      \textbf{3} the device reads the data from the buffer}};
  \end{scope}
\end{tikzpicture}

  \caption{Sending a packet.  (a) With programmed I/O the CPU copies the
    packet into the data register piece by piece.  (b) With DMA the CPU only
    writes the command and configures DMA; the device reads the packet from
    memory.}
  \label{fig:l11-pio-dma}
\end{figure}

\begin{example}
  A network interface with an 8-byte data register sends a 9000-byte frame.
  With PIO the CPU performs 1 access for the command and $9000/8 = 1125$
  accesses for the data, 1126 accesses in total.  With DMA it performs 2: one
  for the command and one for the DMA configuration.
\end{example}

DMA needs fewer steps, but configuring DMA is a more complex operation than
a single register write.  For small transfers, a few register writes with
PIO can therefore cost less than one DMA configuration: PIO is cheaper
whenever the data accesses it needs cost less than the configuration of
DMA.\tip{DMA pays off when a transfer spans many register widths; for
transfers of a few bytes PIO is cheaper.}  A further requirement arises
because the device accesses the buffer by its physical address: the buffer must remain in physical memory until
the transfer completes.  Its pages are therefore pinned
(Lesson~8)\index{pinned page}, that is, excluded from swapping; otherwise
the device could read or overwrite whatever page had replaced the buffer.

\begin{keypoint}
  The CPU reaches device registers through memory-mapped I/O or I/O ports;
  a device reaches the CPU through interrupts, or the operating system polls
  it.  Data move either through the CPU register by register (PIO), or
  directly between a memory buffer and the device (DMA), which saves CPU
  accesses at the price of a more complex configuration and pinned memory.
\end{keypoint}

\section{Typical device access}
\label{sec:l11-access}

An application does not normally interact with a device by itself; its
request passes through several layers of the operating
system.\source{P3L5, typical device access, OS bypass; OSTEP ch.~36.}  The
typical path of a device access is as follows (\cref{fig:l11-access}a):
\begin{enumerate}
  \item The process performs a system call that requests an operation, for
    example to send data over a socket or to read from a file.
  \item The kernel executes the in-kernel stack associated with the device:
    the TCP/IP stack forms a packet, or the file system determines which
    block holds the requested data.
  \item The kernel invokes the appropriate device driver, here the network
    driver or the block device driver.
  \item The driver configures the request on the device, with PIO register
    writes or a DMA configuration.
  \item The device performs the request: it transmits the packet or reads
    the blocks from the disk.
\end{enumerate}
Results and completion notifications travel back the same way: the device
raises an interrupt or is polled, the driver passes the data up the stack,
and the system call returns to the process.

\begin{figure}[htbp]
  \centering
  % (a) Typical device access through the kernel; (b) operating system bypass
% with a user-level driver.
% Usage: % (a) Typical device access through the kernel; (b) operating system bypass
% with a user-level driver.
% Usage: % (a) Typical device access through the kernel; (b) operating system bypass
% with a user-level driver.
% Usage: \input{figures/l11-access}   (width ~118 mm)
\begin{tikzpicture}[x=1mm, y=1mm, font=\scriptsize\sffamily,
  >={Stealth[length=1.5mm]},
  box/.style={draw, fill=white, minimum height=7mm, inner sep=1pt, align=center},
  kbox/.style={draw, fill=ln-box, minimum height=7mm, inner sep=1pt, align=center},
  dev/.style={draw, fill=ln-accent!22, minimum height=7mm, inner sep=1pt, align=center},
  lab/.style={font=\scriptsize\sffamily, align=center},
  side/.style={font=\scriptsize\sffamily\itshape, text=ln-muted},
  ttl/.style={font=\footnotesize\sffamily\bfseries, anchor=west}]
  % boundaries
  \draw[densely dashed, ln-muted] (0,23) -- (118,23);
  \draw[densely dashed, ln-muted] (0,-11) -- (118,-11);
  \node[side, anchor=south west] at (0,23.3) {\iflangja{ユーザ}{user}};
  \node[side, anchor=north west] at (0,22.7) {\iflangja{カーネル}{kernel}};
  \node[side, anchor=north west] at (0,-11.3) {\iflangja{デバイス}{device}};
  % ---------------- (a) typical access
  \node[ttl] at (12,40) {\iflangja{(a) カーネルを経由するアクセス}{(a) access through the kernel}};
  \node[box, minimum width=34mm] (p) at (34,31) {\iflangja{ユーザプロセス}{user process}};
  \node[kbox, minimum width=34mm] (st) at (34,11)
    {\iflangja{カーネル内スタック\\（TCP/IP，ファイルシステム）}{in-kernel stack\\(TCP/IP, file system)}};
  \node[kbox, minimum width=34mm] (dr) at (34,-2) {\iflangja{デバイスドライバ}{device driver}};
  \node[dev, minimum width=34mm] (d) at (34,-21) {\iflangja{デバイス}{device}};
  \draw[->, thick] ([xshift=-6mm]p.south) -- node[lab, left, pos=0.72] {\iflangja{システムコール}{system call}} ([xshift=-6mm]st.north);
  \draw[->, thick] ([xshift=-6mm]st.south) -- ([xshift=-6mm]dr.north);
  \draw[->, thick] ([xshift=-6mm]dr.south) -- node[lab, left, pos=0.22] {PIO/DMA} ([xshift=-6mm]d.north);
  \draw[->] ([xshift=6mm]d.north) -- node[lab, right, pos=0.78] {\iflangja{割り込み}{interrupt}} ([xshift=6mm]dr.south);
  \draw[->] ([xshift=6mm]dr.north) -- ([xshift=6mm]st.south);
  \draw[->] ([xshift=6mm]st.north) -- node[lab, right, pos=0.28] {\iflangja{結果}{result}} ([xshift=6mm]p.south);
  % ---------------- (b) OS bypass
  \node[ttl] at (66,40) {\iflangja{(b) OS バイパス}{(b) operating system bypass}};
  \draw[fill=white] (66,25.5) rectangle (100,37);
  \node[lab] at (83,34.5) {\iflangja{ユーザプロセス}{user process}};
  \node[kbox, minimum width=30mm, minimum height=5mm] (ud) at (83,29.5)
    {\iflangja{ユーザレベルドライバ}{user-level driver}};
  \node[kbox, minimum width=26mm] (kd) at (104,5)
    {\iflangja{カーネルの\\ドライバ}{kernel driver}};
  \node[dev, minimum width=46mm] (d2) at (90,-21) {\iflangja{デバイス}{device}};
  \draw[<->, line width=1.4pt, ln-accent] (77,27) -- node[lab, fill=white, text=black, inner sep=0.8pt]
    {\iflangja{レジスタと\\メモリに\\直接アクセス}{direct access\\to registers\\and memory}} (77,-17.5);
  \draw[->, densely dashed] (kd.south) -- node[lab, right, pos=0.3] {\iflangja{設定}{configure}} (kd.south |- d2.north);
  \draw[->, densely dashed] (kd.north) |- node[lab, pos=0.15, right] {\iflangja{許可}{grant}} (100,31);
\end{tikzpicture}
   (width ~118 mm)
\begin{tikzpicture}[x=1mm, y=1mm, font=\scriptsize\sffamily,
  >={Stealth[length=1.5mm]},
  box/.style={draw, fill=white, minimum height=7mm, inner sep=1pt, align=center},
  kbox/.style={draw, fill=ln-box, minimum height=7mm, inner sep=1pt, align=center},
  dev/.style={draw, fill=ln-accent!22, minimum height=7mm, inner sep=1pt, align=center},
  lab/.style={font=\scriptsize\sffamily, align=center},
  side/.style={font=\scriptsize\sffamily\itshape, text=ln-muted},
  ttl/.style={font=\footnotesize\sffamily\bfseries, anchor=west}]
  % boundaries
  \draw[densely dashed, ln-muted] (0,23) -- (118,23);
  \draw[densely dashed, ln-muted] (0,-11) -- (118,-11);
  \node[side, anchor=south west] at (0,23.3) {\iflangja{ユーザ}{user}};
  \node[side, anchor=north west] at (0,22.7) {\iflangja{カーネル}{kernel}};
  \node[side, anchor=north west] at (0,-11.3) {\iflangja{デバイス}{device}};
  % ---------------- (a) typical access
  \node[ttl] at (12,40) {\iflangja{(a) カーネルを経由するアクセス}{(a) access through the kernel}};
  \node[box, minimum width=34mm] (p) at (34,31) {\iflangja{ユーザプロセス}{user process}};
  \node[kbox, minimum width=34mm] (st) at (34,11)
    {\iflangja{カーネル内スタック\\（TCP/IP，ファイルシステム）}{in-kernel stack\\(TCP/IP, file system)}};
  \node[kbox, minimum width=34mm] (dr) at (34,-2) {\iflangja{デバイスドライバ}{device driver}};
  \node[dev, minimum width=34mm] (d) at (34,-21) {\iflangja{デバイス}{device}};
  \draw[->, thick] ([xshift=-6mm]p.south) -- node[lab, left, pos=0.72] {\iflangja{システムコール}{system call}} ([xshift=-6mm]st.north);
  \draw[->, thick] ([xshift=-6mm]st.south) -- ([xshift=-6mm]dr.north);
  \draw[->, thick] ([xshift=-6mm]dr.south) -- node[lab, left, pos=0.22] {PIO/DMA} ([xshift=-6mm]d.north);
  \draw[->] ([xshift=6mm]d.north) -- node[lab, right, pos=0.78] {\iflangja{割り込み}{interrupt}} ([xshift=6mm]dr.south);
  \draw[->] ([xshift=6mm]dr.north) -- ([xshift=6mm]st.south);
  \draw[->] ([xshift=6mm]st.north) -- node[lab, right, pos=0.28] {\iflangja{結果}{result}} ([xshift=6mm]p.south);
  % ---------------- (b) OS bypass
  \node[ttl] at (66,40) {\iflangja{(b) OS バイパス}{(b) operating system bypass}};
  \draw[fill=white] (66,25.5) rectangle (100,37);
  \node[lab] at (83,34.5) {\iflangja{ユーザプロセス}{user process}};
  \node[kbox, minimum width=30mm, minimum height=5mm] (ud) at (83,29.5)
    {\iflangja{ユーザレベルドライバ}{user-level driver}};
  \node[kbox, minimum width=26mm] (kd) at (104,5)
    {\iflangja{カーネルの\\ドライバ}{kernel driver}};
  \node[dev, minimum width=46mm] (d2) at (90,-21) {\iflangja{デバイス}{device}};
  \draw[<->, line width=1.4pt, ln-accent] (77,27) -- node[lab, fill=white, text=black, inner sep=0.8pt]
    {\iflangja{レジスタと\\メモリに\\直接アクセス}{direct access\\to registers\\and memory}} (77,-17.5);
  \draw[->, densely dashed] (kd.south) -- node[lab, right, pos=0.3] {\iflangja{設定}{configure}} (kd.south |- d2.north);
  \draw[->, densely dashed] (kd.north) |- node[lab, pos=0.15, right] {\iflangja{許可}{grant}} (100,31);
\end{tikzpicture}
   (width ~118 mm)
\begin{tikzpicture}[x=1mm, y=1mm, font=\scriptsize\sffamily,
  >={Stealth[length=1.5mm]},
  box/.style={draw, fill=white, minimum height=7mm, inner sep=1pt, align=center},
  kbox/.style={draw, fill=ln-box, minimum height=7mm, inner sep=1pt, align=center},
  dev/.style={draw, fill=ln-accent!22, minimum height=7mm, inner sep=1pt, align=center},
  lab/.style={font=\scriptsize\sffamily, align=center},
  side/.style={font=\scriptsize\sffamily\itshape, text=ln-muted},
  ttl/.style={font=\footnotesize\sffamily\bfseries, anchor=west}]
  % boundaries
  \draw[densely dashed, ln-muted] (0,23) -- (118,23);
  \draw[densely dashed, ln-muted] (0,-11) -- (118,-11);
  \node[side, anchor=south west] at (0,23.3) {\iflangja{ユーザ}{user}};
  \node[side, anchor=north west] at (0,22.7) {\iflangja{カーネル}{kernel}};
  \node[side, anchor=north west] at (0,-11.3) {\iflangja{デバイス}{device}};
  % ---------------- (a) typical access
  \node[ttl] at (12,40) {\iflangja{(a) カーネルを経由するアクセス}{(a) access through the kernel}};
  \node[box, minimum width=34mm] (p) at (34,31) {\iflangja{ユーザプロセス}{user process}};
  \node[kbox, minimum width=34mm] (st) at (34,11)
    {\iflangja{カーネル内スタック\\（TCP/IP，ファイルシステム）}{in-kernel stack\\(TCP/IP, file system)}};
  \node[kbox, minimum width=34mm] (dr) at (34,-2) {\iflangja{デバイスドライバ}{device driver}};
  \node[dev, minimum width=34mm] (d) at (34,-21) {\iflangja{デバイス}{device}};
  \draw[->, thick] ([xshift=-6mm]p.south) -- node[lab, left, pos=0.72] {\iflangja{システムコール}{system call}} ([xshift=-6mm]st.north);
  \draw[->, thick] ([xshift=-6mm]st.south) -- ([xshift=-6mm]dr.north);
  \draw[->, thick] ([xshift=-6mm]dr.south) -- node[lab, left, pos=0.22] {PIO/DMA} ([xshift=-6mm]d.north);
  \draw[->] ([xshift=6mm]d.north) -- node[lab, right, pos=0.78] {\iflangja{割り込み}{interrupt}} ([xshift=6mm]dr.south);
  \draw[->] ([xshift=6mm]dr.north) -- ([xshift=6mm]st.south);
  \draw[->] ([xshift=6mm]st.north) -- node[lab, right, pos=0.28] {\iflangja{結果}{result}} ([xshift=6mm]p.south);
  % ---------------- (b) OS bypass
  \node[ttl] at (66,40) {\iflangja{(b) OS バイパス}{(b) operating system bypass}};
  \draw[fill=white] (66,25.5) rectangle (100,37);
  \node[lab] at (83,34.5) {\iflangja{ユーザプロセス}{user process}};
  \node[kbox, minimum width=30mm, minimum height=5mm] (ud) at (83,29.5)
    {\iflangja{ユーザレベルドライバ}{user-level driver}};
  \node[kbox, minimum width=26mm] (kd) at (104,5)
    {\iflangja{カーネルの\\ドライバ}{kernel driver}};
  \node[dev, minimum width=46mm] (d2) at (90,-21) {\iflangja{デバイス}{device}};
  \draw[<->, line width=1.4pt, ln-accent] (77,27) -- node[lab, fill=white, text=black, inner sep=0.8pt]
    {\iflangja{レジスタと\\メモリに\\直接アクセス}{direct access\\to registers\\and memory}} (77,-17.5);
  \draw[->, densely dashed] (kd.south) -- node[lab, right, pos=0.3] {\iflangja{設定}{configure}} (kd.south |- d2.north);
  \draw[->, densely dashed] (kd.north) |- node[lab, pos=0.15, right] {\iflangja{許可}{grant}} (100,31);
\end{tikzpicture}

  \caption{(a) A typical device access passes through the system call, the
    in-kernel stack and the device driver.  (b) With operating system bypass
    the process accesses the device through a user-level driver; the kernel
    only configures and controls the access.}
  \label{fig:l11-access}
\end{figure}

\subsection{Operating system bypass}

Going through the kernel is not strictly necessary.

\begin{definition}[Operating system bypass]
  A method of device access in which the registers and memory of a device
  are made directly accessible to a user process is called
  \term{operating system bypass}[OSバイパス].  The operating system configures the access and then stays
  out of the way, so that the process performs device operations without
  executing kernel code (\cref{fig:l11-access}b).
\end{definition}

The code that would otherwise run in the kernel driver is then provided by a
\emph{user-level driver}, a library linked into the application.  The
operating system retains coarse-grained control: it can, for example,
enable or disable the device, or grant additional processes access to it.
Bypass relies on features of the device:
\begin{itemize}
  \item The device must have sufficient registers, so that some of them can
    be mapped into a user process while those needed for control remain
    accessible only to the kernel.
  \item If several processes use the device, the device must be able to
    \emph{demultiplex}, that is, determine to which process data belong.  In
    a typical device access the kernel does this: its network stack, for
    example, inspects the headers of an arriving packet to find the
    receiving process.  With bypass the device itself must do so.
\end{itemize}

\section{Synchronous and asynchronous access}
\label{sec:l11-sync}

While a device carries out a request, the thread that issued it can either
wait or continue.\source{P3L5, sync vs.\ async access; Silberschatz et al.\
ch.~12.}

\begin{definition}[Synchronous I/O]
  With \term{synchronous I/O}[同期I/O] a thread that issues an I/O request
  is blocked until the request completes.  The operating system places the
  thread on the wait queue associated with the device, the I/O queue of
  Lesson~2, and makes it runnable again when the device reports completion.
\end{definition}

With \term*{asynchronous I/O}\index{asynchronous I/O} (Lesson~6) the call
returns at once and the thread continues executing
(\cref{fig:l11-sync-async}).  It learns the outcome later in one of two
ways: either the thread itself checks whether the operation has completed
and then retrieves the results, or the operating system notifies it that
the operation has completed and that the results are available.  When the
kernel does not support asynchronous access to a device, it can be emulated
in user space with helper threads or processes that issue the blocking
calls on behalf of the main thread, as the Flash web server of Lesson~6
does.

\begin{figure}[htbp]
  \centering
  % Synchronous versus asynchronous I/O: timelines of the requesting thread and
% the device.
% Usage: % Synchronous versus asynchronous I/O: timelines of the requesting thread and
% the device.
% Usage: % Synchronous versus asynchronous I/O: timelines of the requesting thread and
% the device.
% Usage: \input{figures/l11-sync-async}   (width ~116 mm)
\begin{tikzpicture}[x=1mm, y=1mm, font=\footnotesize\sffamily,
  >={Stealth[length=1.6mm]},
  lab/.style={font=\scriptsize\sffamily, align=center},
  row/.style={font=\scriptsize\sffamily, anchor=east},
  ttl/.style={font=\footnotesize\sffamily\bfseries, anchor=west}]
  % \lelseg{x1}{x2}{y}{fill}{text}
  \def\lelseg#1#2#3#4#5{%
    \draw[fill=#4] (#1,#3-2.5) rectangle (#2,#3+2.5);
    \node[lab] at ({(#1+#2)/2},#3) {#5};}
  % (a) synchronous
  \node[ttl] at (0,31) {\iflangja{(a) 同期I/O}{(a) synchronous I/O}};
  \node[row] at (20,24) {\iflangja{スレッド}{thread}};
  \lelseg{22}{44}{24}{white}{\iflangja{実行}{running}}
  \lelseg{44}{86}{24}{black!15}{\iflangja{ブロック（デバイスの待ち行列）}{blocked (on the device's wait queue)}}
  \lelseg{86}{116}{24}{white}{\iflangja{結果を使って続行}{continues with result}}
  \node[row] at (20,15) {\iflangja{デバイス}{device}};
  \lelseg{48}{80}{15}{ln-box}{\iflangja{I/O 処理}{I/O operation}}
  \draw[->] (44,21.5) -- node[lab, left=0.3mm, pos=0.35] {\iflangja{要求}{request}} (48,17.5);
  \draw[->, ln-caution] (80,17.5) -- node[lab, right=0.3mm, pos=0.55, text=ln-caution] {\iflangja{完了}{done}} (86,21.5);
  % (b) asynchronous
  \node[ttl] at (0,6) {\iflangja{(b) 非同期I/O}{(b) asynchronous I/O}};
  \node[row] at (20,-1) {\iflangja{スレッド}{thread}};
  \lelseg{22}{44}{-1}{white}{\iflangja{実行}{running}}
  \lelseg{44}{92}{-1}{white}{\iflangja{他の処理を続ける}{continues with other work}}
  \lelseg{92}{116}{-1}{ln-accent!22}{\iflangja{結果を取得}{retrieves result}}
  \node[row] at (20,-10) {\iflangja{デバイス}{device}};
  \lelseg{48}{80}{-10}{ln-box}{\iflangja{I/O 処理}{I/O operation}}
  \draw[->] (44,-3.5) -- node[lab, left=0.3mm, pos=0.35] {\iflangja{要求}{request}} (48,-7.5);
  \draw[->, densely dashed, ln-caution] (80,-7.5) -- (92,-3.5);
  \node[lab, anchor=north west, text=ln-caution] at (81,-12.8)
    {\iflangja{通知，またはスレッドが確認}{notification, or the thread checks}};
  % time axis
  \draw[->] (22,-21) -- node[lab, below] {\iflangja{時間}{time}} (116,-21);
\end{tikzpicture}
   (width ~116 mm)
\begin{tikzpicture}[x=1mm, y=1mm, font=\footnotesize\sffamily,
  >={Stealth[length=1.6mm]},
  lab/.style={font=\scriptsize\sffamily, align=center},
  row/.style={font=\scriptsize\sffamily, anchor=east},
  ttl/.style={font=\footnotesize\sffamily\bfseries, anchor=west}]
  % \lelseg{x1}{x2}{y}{fill}{text}
  \def\lelseg#1#2#3#4#5{%
    \draw[fill=#4] (#1,#3-2.5) rectangle (#2,#3+2.5);
    \node[lab] at ({(#1+#2)/2},#3) {#5};}
  % (a) synchronous
  \node[ttl] at (0,31) {\iflangja{(a) 同期I/O}{(a) synchronous I/O}};
  \node[row] at (20,24) {\iflangja{スレッド}{thread}};
  \lelseg{22}{44}{24}{white}{\iflangja{実行}{running}}
  \lelseg{44}{86}{24}{black!15}{\iflangja{ブロック（デバイスの待ち行列）}{blocked (on the device's wait queue)}}
  \lelseg{86}{116}{24}{white}{\iflangja{結果を使って続行}{continues with result}}
  \node[row] at (20,15) {\iflangja{デバイス}{device}};
  \lelseg{48}{80}{15}{ln-box}{\iflangja{I/O 処理}{I/O operation}}
  \draw[->] (44,21.5) -- node[lab, left=0.3mm, pos=0.35] {\iflangja{要求}{request}} (48,17.5);
  \draw[->, ln-caution] (80,17.5) -- node[lab, right=0.3mm, pos=0.55, text=ln-caution] {\iflangja{完了}{done}} (86,21.5);
  % (b) asynchronous
  \node[ttl] at (0,6) {\iflangja{(b) 非同期I/O}{(b) asynchronous I/O}};
  \node[row] at (20,-1) {\iflangja{スレッド}{thread}};
  \lelseg{22}{44}{-1}{white}{\iflangja{実行}{running}}
  \lelseg{44}{92}{-1}{white}{\iflangja{他の処理を続ける}{continues with other work}}
  \lelseg{92}{116}{-1}{ln-accent!22}{\iflangja{結果を取得}{retrieves result}}
  \node[row] at (20,-10) {\iflangja{デバイス}{device}};
  \lelseg{48}{80}{-10}{ln-box}{\iflangja{I/O 処理}{I/O operation}}
  \draw[->] (44,-3.5) -- node[lab, left=0.3mm, pos=0.35] {\iflangja{要求}{request}} (48,-7.5);
  \draw[->, densely dashed, ln-caution] (80,-7.5) -- (92,-3.5);
  \node[lab, anchor=north west, text=ln-caution] at (81,-12.8)
    {\iflangja{通知，またはスレッドが確認}{notification, or the thread checks}};
  % time axis
  \draw[->] (22,-21) -- node[lab, below] {\iflangja{時間}{time}} (116,-21);
\end{tikzpicture}
   (width ~116 mm)
\begin{tikzpicture}[x=1mm, y=1mm, font=\footnotesize\sffamily,
  >={Stealth[length=1.6mm]},
  lab/.style={font=\scriptsize\sffamily, align=center},
  row/.style={font=\scriptsize\sffamily, anchor=east},
  ttl/.style={font=\footnotesize\sffamily\bfseries, anchor=west}]
  % \lelseg{x1}{x2}{y}{fill}{text}
  \def\lelseg#1#2#3#4#5{%
    \draw[fill=#4] (#1,#3-2.5) rectangle (#2,#3+2.5);
    \node[lab] at ({(#1+#2)/2},#3) {#5};}
  % (a) synchronous
  \node[ttl] at (0,31) {\iflangja{(a) 同期I/O}{(a) synchronous I/O}};
  \node[row] at (20,24) {\iflangja{スレッド}{thread}};
  \lelseg{22}{44}{24}{white}{\iflangja{実行}{running}}
  \lelseg{44}{86}{24}{black!15}{\iflangja{ブロック（デバイスの待ち行列）}{blocked (on the device's wait queue)}}
  \lelseg{86}{116}{24}{white}{\iflangja{結果を使って続行}{continues with result}}
  \node[row] at (20,15) {\iflangja{デバイス}{device}};
  \lelseg{48}{80}{15}{ln-box}{\iflangja{I/O 処理}{I/O operation}}
  \draw[->] (44,21.5) -- node[lab, left=0.3mm, pos=0.35] {\iflangja{要求}{request}} (48,17.5);
  \draw[->, ln-caution] (80,17.5) -- node[lab, right=0.3mm, pos=0.55, text=ln-caution] {\iflangja{完了}{done}} (86,21.5);
  % (b) asynchronous
  \node[ttl] at (0,6) {\iflangja{(b) 非同期I/O}{(b) asynchronous I/O}};
  \node[row] at (20,-1) {\iflangja{スレッド}{thread}};
  \lelseg{22}{44}{-1}{white}{\iflangja{実行}{running}}
  \lelseg{44}{92}{-1}{white}{\iflangja{他の処理を続ける}{continues with other work}}
  \lelseg{92}{116}{-1}{ln-accent!22}{\iflangja{結果を取得}{retrieves result}}
  \node[row] at (20,-10) {\iflangja{デバイス}{device}};
  \lelseg{48}{80}{-10}{ln-box}{\iflangja{I/O 処理}{I/O operation}}
  \draw[->] (44,-3.5) -- node[lab, left=0.3mm, pos=0.35] {\iflangja{要求}{request}} (48,-7.5);
  \draw[->, densely dashed, ln-caution] (80,-7.5) -- (92,-3.5);
  \node[lab, anchor=north west, text=ln-caution] at (81,-12.8)
    {\iflangja{通知，またはスレッドが確認}{notification, or the thread checks}};
  % time axis
  \draw[->] (22,-21) -- node[lab, below] {\iflangja{時間}{time}} (116,-21);
\end{tikzpicture}

  \caption{(a) With synchronous I/O the thread blocks until the device
    completes the operation.  (b) With asynchronous I/O the thread continues
    and retrieves the result after being notified or after checking.}
  \label{fig:l11-sync-async}
\end{figure}

\begin{keypoint}
  A device access normally passes from the system call through the
  in-kernel stack and the driver to the device, and the result returns the
  same way.  Operating system bypass removes the kernel from this path when
  the device can protect and demultiplex its own resources.  Independently
  of the path, the requesting thread either blocks (synchronous I/O) or
  continues and collects the result later (asynchronous I/O).
\end{keypoint}

\section{Block device stack}
\label{sec:l11-block-stack}

Processes rarely address a block device by block numbers.\source{P3L5,
block device stack; OSTEP ch.~36, 39.}  They use files, the logical storage
units, through the POSIX API with operations such as \texttt{open},
\texttt{read}, \texttt{write} and \texttt{close}.
Beneath the API the kernel is organized in layers, each of which uses the
interface of the layer below it (\cref{fig:l11-block-stack}).

\begin{definition}[File system]
  A \term{file system}[ファイルシステム] is the kernel component that
  organizes the blocks of a storage device into files and directories: it
  determines where the data of each file reside, and how a file is found and
  accessed.
\end{definition}

The operating system specifies the interface that every file system
implements.  Below the file systems lies a layer that hides the differences
between block devices.

\begin{definition}[Generic block layer]
  The \term{generic block layer}[汎用ブロック層] is the kernel layer that
  provides a single, operating-system-standardized interface for reading
  and writing blocks on any block device, and that translates each request
  into the interface of the driver of the particular device.
\end{definition}

The same file system therefore works on every kind of block device: the
generic block layer adapts each request to the driver of the device, and the
driver speaks the protocol of its device, such as SCSI or ATA, towards the
hardware.

\begin{figure}[htbp]
  \centering
  % The block device stack: file system, generic block layer and device driver
% between a process and a block device.
% Usage: % The block device stack: file system, generic block layer and device driver
% between a process and a block device.
% Usage: % The block device stack: file system, generic block layer and device driver
% between a process and a block device.
% Usage: \input{figures/l11-block-stack}   (width ~112 mm)
\begin{tikzpicture}[x=1mm, y=1mm, font=\footnotesize\sffamily,
  >={Stealth[length=1.6mm]},
  box/.style={draw, fill=white, minimum width=38mm, minimum height=7mm, inner sep=1pt},
  kbox/.style={draw, fill=ln-box, minimum width=38mm, minimum height=7mm, inner sep=1pt},
  dev/.style={draw, fill=ln-accent!22, minimum width=38mm, minimum height=7mm, inner sep=1pt},
  ifc/.style={font=\scriptsize\sffamily, anchor=west, align=left},
  side/.style={font=\scriptsize\sffamily\itshape, text=ln-muted, anchor=west}]
  \node[box]  (p)  at (34,40) {\iflangja{ユーザプロセス}{user process}};
  \node[kbox] (fs) at (34,27) {\iflangja{ファイルシステム}{file system}};
  \node[kbox] (gb) at (34,14) {\iflangja{汎用ブロック層}{generic block layer}};
  \node[kbox] (dd) at (34,1)  {\iflangja{デバイスドライバ}{device driver}};
  \node[dev]  (bd) at (34,-12) {\iflangja{ブロックデバイス}{block device}};
  \foreach \a/\b in {p/fs, fs/gb, gb/dd, dd/bd} \draw[<->, thick] (\a.south) -- (\b.north);
  \node[ifc] at (40,33.5) {\iflangja{POSIX API：open, read, write, close}{POSIX API: open, read, write, close}};
  \node[ifc] at (40,20.5) {\iflangja{汎用ブロックインタフェース：ブロックの読み書き}{generic block interface: read or write blocks}};
  \node[ifc] at (40,7.5)  {\iflangja{プロトコル固有のインタフェース}{protocol-specific interface}};
  \node[ifc] at (40,-5.5) {\iflangja{デバイスのプロトコル（SCSI, ATA など）}{device protocol (e.g.\ SCSI, ATA)}};
  % boundaries
  \draw[densely dashed, ln-muted] (0,33.5) -- (15,33.5);
  \draw[densely dashed, ln-muted] (0,-5.5) -- (15,-5.5);
  \node[side, anchor=south west] at (0,33.8) {\iflangja{ユーザ}{user}};
  \node[side, anchor=north west] at (0,33.2) {\iflangja{カーネル}{kernel}};
  \node[side, anchor=north west] at (0,-5.8) {\iflangja{ハードウェア}{hardware}};
  % ioctl
  \draw[->, densely dashed, ln-accent] (p.west) .. controls (10,34) and (10,7) .. (dd.west);
  \node[font=\scriptsize\ttfamily, text=ln-accent, fill=white, inner sep=0.8pt] at (11.2,20.5) {ioctl};
\end{tikzpicture}
   (width ~112 mm)
\begin{tikzpicture}[x=1mm, y=1mm, font=\footnotesize\sffamily,
  >={Stealth[length=1.6mm]},
  box/.style={draw, fill=white, minimum width=38mm, minimum height=7mm, inner sep=1pt},
  kbox/.style={draw, fill=ln-box, minimum width=38mm, minimum height=7mm, inner sep=1pt},
  dev/.style={draw, fill=ln-accent!22, minimum width=38mm, minimum height=7mm, inner sep=1pt},
  ifc/.style={font=\scriptsize\sffamily, anchor=west, align=left},
  side/.style={font=\scriptsize\sffamily\itshape, text=ln-muted, anchor=west}]
  \node[box]  (p)  at (34,40) {\iflangja{ユーザプロセス}{user process}};
  \node[kbox] (fs) at (34,27) {\iflangja{ファイルシステム}{file system}};
  \node[kbox] (gb) at (34,14) {\iflangja{汎用ブロック層}{generic block layer}};
  \node[kbox] (dd) at (34,1)  {\iflangja{デバイスドライバ}{device driver}};
  \node[dev]  (bd) at (34,-12) {\iflangja{ブロックデバイス}{block device}};
  \foreach \a/\b in {p/fs, fs/gb, gb/dd, dd/bd} \draw[<->, thick] (\a.south) -- (\b.north);
  \node[ifc] at (40,33.5) {\iflangja{POSIX API：open, read, write, close}{POSIX API: open, read, write, close}};
  \node[ifc] at (40,20.5) {\iflangja{汎用ブロックインタフェース：ブロックの読み書き}{generic block interface: read or write blocks}};
  \node[ifc] at (40,7.5)  {\iflangja{プロトコル固有のインタフェース}{protocol-specific interface}};
  \node[ifc] at (40,-5.5) {\iflangja{デバイスのプロトコル（SCSI, ATA など）}{device protocol (e.g.\ SCSI, ATA)}};
  % boundaries
  \draw[densely dashed, ln-muted] (0,33.5) -- (15,33.5);
  \draw[densely dashed, ln-muted] (0,-5.5) -- (15,-5.5);
  \node[side, anchor=south west] at (0,33.8) {\iflangja{ユーザ}{user}};
  \node[side, anchor=north west] at (0,33.2) {\iflangja{カーネル}{kernel}};
  \node[side, anchor=north west] at (0,-5.8) {\iflangja{ハードウェア}{hardware}};
  % ioctl
  \draw[->, densely dashed, ln-accent] (p.west) .. controls (10,34) and (10,7) .. (dd.west);
  \node[font=\scriptsize\ttfamily, text=ln-accent, fill=white, inner sep=0.8pt] at (11.2,20.5) {ioctl};
\end{tikzpicture}
   (width ~112 mm)
\begin{tikzpicture}[x=1mm, y=1mm, font=\footnotesize\sffamily,
  >={Stealth[length=1.6mm]},
  box/.style={draw, fill=white, minimum width=38mm, minimum height=7mm, inner sep=1pt},
  kbox/.style={draw, fill=ln-box, minimum width=38mm, minimum height=7mm, inner sep=1pt},
  dev/.style={draw, fill=ln-accent!22, minimum width=38mm, minimum height=7mm, inner sep=1pt},
  ifc/.style={font=\scriptsize\sffamily, anchor=west, align=left},
  side/.style={font=\scriptsize\sffamily\itshape, text=ln-muted, anchor=west}]
  \node[box]  (p)  at (34,40) {\iflangja{ユーザプロセス}{user process}};
  \node[kbox] (fs) at (34,27) {\iflangja{ファイルシステム}{file system}};
  \node[kbox] (gb) at (34,14) {\iflangja{汎用ブロック層}{generic block layer}};
  \node[kbox] (dd) at (34,1)  {\iflangja{デバイスドライバ}{device driver}};
  \node[dev]  (bd) at (34,-12) {\iflangja{ブロックデバイス}{block device}};
  \foreach \a/\b in {p/fs, fs/gb, gb/dd, dd/bd} \draw[<->, thick] (\a.south) -- (\b.north);
  \node[ifc] at (40,33.5) {\iflangja{POSIX API：open, read, write, close}{POSIX API: open, read, write, close}};
  \node[ifc] at (40,20.5) {\iflangja{汎用ブロックインタフェース：ブロックの読み書き}{generic block interface: read or write blocks}};
  \node[ifc] at (40,7.5)  {\iflangja{プロトコル固有のインタフェース}{protocol-specific interface}};
  \node[ifc] at (40,-5.5) {\iflangja{デバイスのプロトコル（SCSI, ATA など）}{device protocol (e.g.\ SCSI, ATA)}};
  % boundaries
  \draw[densely dashed, ln-muted] (0,33.5) -- (15,33.5);
  \draw[densely dashed, ln-muted] (0,-5.5) -- (15,-5.5);
  \node[side, anchor=south west] at (0,33.8) {\iflangja{ユーザ}{user}};
  \node[side, anchor=north west] at (0,33.2) {\iflangja{カーネル}{kernel}};
  \node[side, anchor=north west] at (0,-5.8) {\iflangja{ハードウェア}{hardware}};
  % ioctl
  \draw[->, densely dashed, ln-accent] (p.west) .. controls (10,34) and (10,7) .. (dd.west);
  \node[font=\scriptsize\ttfamily, text=ln-accent, fill=white, inner sep=0.8pt] at (11.2,20.5) {ioctl};
\end{tikzpicture}

  \caption{The block device stack.  Each layer uses the interface shown
    below it; \texttt{ioctl} carries device-specific requests to the
    driver.}
  \label{fig:l11-block-stack}
\end{figure}

Some operations on a device fit neither \texttt{read} nor \texttt{write}.
Linux and other UNIX-like systems provide the system call
\term[ioctl]{\texttt{ioctl}}[ioctl] (\enquote{I/O control}) through which a
process accesses and manipulates device-specific information.  It takes a
file descriptor (\cref{sec:l11-vfs-abs}) of an open device file, a request
code defined for the device class or driver, and, depending on the request,
a pointer to request-specific data.

\needspace{11\baselineskip}
\begin{example}
  A program opens the tray of the first optical drive:
  \begin{lstlisting}[language=C, numbers=none]
#include <fcntl.h>
#include <unistd.h>
#include <sys/ioctl.h>
#include <linux/cdrom.h>

int fd = open("/dev/sr0", O_RDONLY | O_NONBLOCK);
if (fd >= 0) {
  ioctl(fd, CDROMEJECT);   /* no read/write equivalent */
  close(fd);
}
\end{lstlisting}
  The flag \texttt{O\_NONBLOCK} allows the device file to be opened even
  when no disc is inserted.
\end{example}

\section{The virtual file system}
\label{sec:l11-vfs}

A single file system on a single local device is not enough for most
systems.\source{P3L5, virtual file system; OSTEP ch.~39; Kerrisk
ch.~14.}  The files may reside on more than one device; different devices
may be served best by different file system implementations; and files may
not reside on a local device at all, but on a remote server that is accessed
over the network (Lesson~14).

\begin{definition}[Virtual file system]
  A \term{virtual file system}[仮想ファイルシステム]%
  \index{VFS|see{virtual file system}} (VFS) is a kernel layer that presents
  one file system interface to applications, and that maps each operation to
  the concrete file system in which the file resides.
\end{definition}

The VFS hides all details of the underlying file systems from applications
(\cref{fig:l11-vfs}).  Applications use the same API regardless of where a
file resides, and the VFS forwards each operation to the file system
concerned, such as a local ext2 file system on a disk or a remote file
system such as NFS.

\begin{figure}[tb]
  \centering
  % The virtual file system layer between the file system interface and the
% concrete local and remote file systems.
% Usage: % The virtual file system layer between the file system interface and the
% concrete local and remote file systems.
% Usage: % The virtual file system layer between the file system interface and the
% concrete local and remote file systems.
% Usage: \input{figures/l11-vfs}   (width ~114 mm)
\begin{tikzpicture}[x=1mm, y=1mm, font=\footnotesize\sffamily,
  >={Stealth[length=1.6mm]},
  bar/.style={draw, fill=ln-box, minimum width=108mm, minimum height=6.5mm, inner sep=1pt},
  fs/.style={draw, fill=white, minimum width=34mm, minimum height=9mm, inner sep=1pt, align=center},
  dev/.style={draw, fill=ln-accent!22, minimum width=22mm, minimum height=6mm, inner sep=1pt}]
  \node[draw, fill=white, minimum width=60mm, minimum height=6.5mm] (app) at (57,36)
    {\iflangja{ユーザプロセス}{user processes}};
  \node[bar] (fsi) at (57,25) {\iflangja{ファイルシステムインタフェース（POSIX API）}{file system interface (POSIX API)}};
  \node[bar, fill=ln-accent!10] (vfs) at (57,14) {\iflangja{仮想ファイルシステム（VFS）}{virtual file system (VFS)}};
  \node[fs] (f1) at (19,0) {\iflangja{ローカル FS\\（例：ext2）}{local file system\\(e.g.\ ext2)}};
  \node[fs] (f2) at (57,0) {\iflangja{ローカル FS\\（別の種類）}{local file system\\(another type)}};
  \node[fs] (f3) at (95,0) {\iflangja{リモート FS\\（例：NFS）}{remote file system\\(e.g.\ NFS)}};
  \node[dev] (d1) at (19,-14) {\iflangja{ディスク}{disk}};
  \node[dev] (d2) at (57,-14) {\iflangja{ディスク}{disk}};
  \node[dev] (d3) at (95,-14) {\iflangja{ネットワーク}{network}};
  \draw[<->, thick] (app) -- (fsi);
  \draw[<->, thick] (fsi) -- (vfs);
  \foreach \x/\f/\d in {19/f1/d1, 57/f2/d2, 95/f3/d3} {
    \draw[<->, thick] (\x,10.75) -- (\f.north);
    \draw[<->, thick] (\f.south) -- (\d.north);
  }
\end{tikzpicture}
   (width ~114 mm)
\begin{tikzpicture}[x=1mm, y=1mm, font=\footnotesize\sffamily,
  >={Stealth[length=1.6mm]},
  bar/.style={draw, fill=ln-box, minimum width=108mm, minimum height=6.5mm, inner sep=1pt},
  fs/.style={draw, fill=white, minimum width=34mm, minimum height=9mm, inner sep=1pt, align=center},
  dev/.style={draw, fill=ln-accent!22, minimum width=22mm, minimum height=6mm, inner sep=1pt}]
  \node[draw, fill=white, minimum width=60mm, minimum height=6.5mm] (app) at (57,36)
    {\iflangja{ユーザプロセス}{user processes}};
  \node[bar] (fsi) at (57,25) {\iflangja{ファイルシステムインタフェース（POSIX API）}{file system interface (POSIX API)}};
  \node[bar, fill=ln-accent!10] (vfs) at (57,14) {\iflangja{仮想ファイルシステム（VFS）}{virtual file system (VFS)}};
  \node[fs] (f1) at (19,0) {\iflangja{ローカル FS\\（例：ext2）}{local file system\\(e.g.\ ext2)}};
  \node[fs] (f2) at (57,0) {\iflangja{ローカル FS\\（別の種類）}{local file system\\(another type)}};
  \node[fs] (f3) at (95,0) {\iflangja{リモート FS\\（例：NFS）}{remote file system\\(e.g.\ NFS)}};
  \node[dev] (d1) at (19,-14) {\iflangja{ディスク}{disk}};
  \node[dev] (d2) at (57,-14) {\iflangja{ディスク}{disk}};
  \node[dev] (d3) at (95,-14) {\iflangja{ネットワーク}{network}};
  \draw[<->, thick] (app) -- (fsi);
  \draw[<->, thick] (fsi) -- (vfs);
  \foreach \x/\f/\d in {19/f1/d1, 57/f2/d2, 95/f3/d3} {
    \draw[<->, thick] (\x,10.75) -- (\f.north);
    \draw[<->, thick] (\f.south) -- (\d.north);
  }
\end{tikzpicture}
   (width ~114 mm)
\begin{tikzpicture}[x=1mm, y=1mm, font=\footnotesize\sffamily,
  >={Stealth[length=1.6mm]},
  bar/.style={draw, fill=ln-box, minimum width=108mm, minimum height=6.5mm, inner sep=1pt},
  fs/.style={draw, fill=white, minimum width=34mm, minimum height=9mm, inner sep=1pt, align=center},
  dev/.style={draw, fill=ln-accent!22, minimum width=22mm, minimum height=6mm, inner sep=1pt}]
  \node[draw, fill=white, minimum width=60mm, minimum height=6.5mm] (app) at (57,36)
    {\iflangja{ユーザプロセス}{user processes}};
  \node[bar] (fsi) at (57,25) {\iflangja{ファイルシステムインタフェース（POSIX API）}{file system interface (POSIX API)}};
  \node[bar, fill=ln-accent!10] (vfs) at (57,14) {\iflangja{仮想ファイルシステム（VFS）}{virtual file system (VFS)}};
  \node[fs] (f1) at (19,0) {\iflangja{ローカル FS\\（例：ext2）}{local file system\\(e.g.\ ext2)}};
  \node[fs] (f2) at (57,0) {\iflangja{ローカル FS\\（別の種類）}{local file system\\(another type)}};
  \node[fs] (f3) at (95,0) {\iflangja{リモート FS\\（例：NFS）}{remote file system\\(e.g.\ NFS)}};
  \node[dev] (d1) at (19,-14) {\iflangja{ディスク}{disk}};
  \node[dev] (d2) at (57,-14) {\iflangja{ディスク}{disk}};
  \node[dev] (d3) at (95,-14) {\iflangja{ネットワーク}{network}};
  \draw[<->, thick] (app) -- (fsi);
  \draw[<->, thick] (fsi) -- (vfs);
  \foreach \x/\f/\d in {19/f1/d1, 57/f2/d2, 95/f3/d3} {
    \draw[<->, thick] (\x,10.75) -- (\f.north);
    \draw[<->, thick] (\f.south) -- (\d.north);
  }
\end{tikzpicture}

  \caption{The virtual file system maps the operations of the file system
    interface to the concrete local and remote file systems.}
  \label{fig:l11-vfs}
\end{figure}

\section{VFS abstractions}
\label{sec:l11-vfs-abs}

Every concrete file system implements a small set of abstractions defined
by the VFS.\source{P3L5, virtual file system abstractions, VFS on disk;
OSTEP ch.~39--40; Kerrisk ch.~14.}  The \emph{file} is the element
on which the VFS operates.  A process refers to an open file through a file
descriptor (\cref{fig:l11-vfs-objects}).

\begin{figure}[tb]
  \centering
  % VFS abstractions: file descriptor, file, dentries and inode in memory;
% superblock, inode blocks and data blocks on disk.
% Usage: % VFS abstractions: file descriptor, file, dentries and inode in memory;
% superblock, inode blocks and data blocks on disk.
% Usage: % VFS abstractions: file descriptor, file, dentries and inode in memory;
% superblock, inode blocks and data blocks on disk.
% Usage: \input{figures/l11-vfs-objects}   (width ~118 mm)
\begin{tikzpicture}[x=1mm, y=1mm, font=\scriptsize\sffamily,
  >={Stealth[length=1.5mm]},
  obj/.style={draw, fill=white, minimum height=8mm, inner sep=1pt, align=center},
  den/.style={draw, fill=white, minimum height=5mm, inner sep=1.2pt, font=\scriptsize\ttfamily},
  lab/.style={font=\scriptsize\sffamily, align=center},
  side/.style={font=\scriptsize\sffamily\bfseries, text=ln-accent}]
  \draw[densely dashed, ln-muted] (0,0) -- (118,0);
  \node[side, anchor=south east] at (118,0.4) {\iflangja{メモリ}{memory}};
  \node[side, anchor=north east] at (118,-0.4) {\iflangja{ディスク}{disk}};
  % file descriptor table of process P
  \node[lab, anchor=south] at (9,42.2) {\iflangja{P の fd}{fds of P}};
  \foreach \i/\y in {0/40,1/36,2/32,3/28} {
    \draw[fill=white] (4,\y-2) rectangle (14,\y+2);
    \node[font=\scriptsize\ttfamily] at (9,\y) {\i};
  }
  \draw[fill=ln-accent!22] (4,26) rectangle (14,30);
  \node[font=\scriptsize\ttfamily] at (9,28) {3};
  % file
  \node[obj, minimum width=26mm] (file) at (36,28)
    {\iflangja{ファイル\\（オフセット，モード）}{file\\(offset, mode)}};
  \draw[->] (14,28) -- (file.west);
  % dentries
  \draw[densely dotted, ln-muted, fill=ln-box] (3,9) rectangle (70,19);
  \node[lab, anchor=north west, text=ln-muted] at (3.5,8.6) {\iflangja{dentry キャッシュ（メモリ上のみ）}{dentry cache (in memory only)}};
  \node[den] (r) at (10,14) {/};
  \node[den] (u) at (24,14) {users};
  \node[den] (j) at (39,14) {joe};
  \node[den] (n) at (58,14) {notes.txt};
  \draw[->] (r) -- (u);  \draw[->] (u) -- (j);  \draw[->] (j) -- (n);
  \draw[->] (file.east) -| (n.north);
  % inode in memory
  \node[obj, minimum width=32mm, minimum height=15mm] (ino) at (97,20)
    {\iflangja{iノード\\所有者，権限，サイズ，\\ブロックの一覧}{inode\\owner, permissions, size,\\list of blocks}};
  \draw[->] (n.east) -- (ino.west |- n.east);
  % on-disk layout
  \def\seg#1#2#3#4{\draw[fill=#3] (#1,-15) rectangle (#2,-8);
    \node[lab] at ({(#1+#2)/2},-11.5) {#4};}
  \seg{4}{30}{ln-accent!22}{\iflangja{スーパーブロック}{superblock}}
  \seg{30}{58}{ln-box}{\iflangja{iノードブロック}{inode blocks}}
  \seg{58}{98}{white}{\iflangja{データブロック}{data blocks}}
  \seg{98}{118}{black!8}{\iflangja{空きブロック}{free blocks}}
  \draw[<->, densely dashed] (85,12.5) -- (48,-8);
  \node[lab, anchor=east] at (64,3.2) {\iflangja{ディスク上の表現}{on-disk copy}};
  \foreach \x in {66,78,90} \draw[->] (97,12.5) -- (\x,-8);
  \draw[->, ln-accent] (17,-8) .. controls (22,-4) and (36,-4) .. (44,-8);
  \node[lab, text=ln-accent, anchor=south] at (30,-3.7) {\iflangja{配置の地図}{map of the layout}};
\end{tikzpicture}
   (width ~118 mm)
\begin{tikzpicture}[x=1mm, y=1mm, font=\scriptsize\sffamily,
  >={Stealth[length=1.5mm]},
  obj/.style={draw, fill=white, minimum height=8mm, inner sep=1pt, align=center},
  den/.style={draw, fill=white, minimum height=5mm, inner sep=1.2pt, font=\scriptsize\ttfamily},
  lab/.style={font=\scriptsize\sffamily, align=center},
  side/.style={font=\scriptsize\sffamily\bfseries, text=ln-accent}]
  \draw[densely dashed, ln-muted] (0,0) -- (118,0);
  \node[side, anchor=south east] at (118,0.4) {\iflangja{メモリ}{memory}};
  \node[side, anchor=north east] at (118,-0.4) {\iflangja{ディスク}{disk}};
  % file descriptor table of process P
  \node[lab, anchor=south] at (9,42.2) {\iflangja{P の fd}{fds of P}};
  \foreach \i/\y in {0/40,1/36,2/32,3/28} {
    \draw[fill=white] (4,\y-2) rectangle (14,\y+2);
    \node[font=\scriptsize\ttfamily] at (9,\y) {\i};
  }
  \draw[fill=ln-accent!22] (4,26) rectangle (14,30);
  \node[font=\scriptsize\ttfamily] at (9,28) {3};
  % file
  \node[obj, minimum width=26mm] (file) at (36,28)
    {\iflangja{ファイル\\（オフセット，モード）}{file\\(offset, mode)}};
  \draw[->] (14,28) -- (file.west);
  % dentries
  \draw[densely dotted, ln-muted, fill=ln-box] (3,9) rectangle (70,19);
  \node[lab, anchor=north west, text=ln-muted] at (3.5,8.6) {\iflangja{dentry キャッシュ（メモリ上のみ）}{dentry cache (in memory only)}};
  \node[den] (r) at (10,14) {/};
  \node[den] (u) at (24,14) {users};
  \node[den] (j) at (39,14) {joe};
  \node[den] (n) at (58,14) {notes.txt};
  \draw[->] (r) -- (u);  \draw[->] (u) -- (j);  \draw[->] (j) -- (n);
  \draw[->] (file.east) -| (n.north);
  % inode in memory
  \node[obj, minimum width=32mm, minimum height=15mm] (ino) at (97,20)
    {\iflangja{iノード\\所有者，権限，サイズ，\\ブロックの一覧}{inode\\owner, permissions, size,\\list of blocks}};
  \draw[->] (n.east) -- (ino.west |- n.east);
  % on-disk layout
  \def\seg#1#2#3#4{\draw[fill=#3] (#1,-15) rectangle (#2,-8);
    \node[lab] at ({(#1+#2)/2},-11.5) {#4};}
  \seg{4}{30}{ln-accent!22}{\iflangja{スーパーブロック}{superblock}}
  \seg{30}{58}{ln-box}{\iflangja{iノードブロック}{inode blocks}}
  \seg{58}{98}{white}{\iflangja{データブロック}{data blocks}}
  \seg{98}{118}{black!8}{\iflangja{空きブロック}{free blocks}}
  \draw[<->, densely dashed] (85,12.5) -- (48,-8);
  \node[lab, anchor=east] at (64,3.2) {\iflangja{ディスク上の表現}{on-disk copy}};
  \foreach \x in {66,78,90} \draw[->] (97,12.5) -- (\x,-8);
  \draw[->, ln-accent] (17,-8) .. controls (22,-4) and (36,-4) .. (44,-8);
  \node[lab, text=ln-accent, anchor=south] at (30,-3.7) {\iflangja{配置の地図}{map of the layout}};
\end{tikzpicture}
   (width ~118 mm)
\begin{tikzpicture}[x=1mm, y=1mm, font=\scriptsize\sffamily,
  >={Stealth[length=1.5mm]},
  obj/.style={draw, fill=white, minimum height=8mm, inner sep=1pt, align=center},
  den/.style={draw, fill=white, minimum height=5mm, inner sep=1.2pt, font=\scriptsize\ttfamily},
  lab/.style={font=\scriptsize\sffamily, align=center},
  side/.style={font=\scriptsize\sffamily\bfseries, text=ln-accent}]
  \draw[densely dashed, ln-muted] (0,0) -- (118,0);
  \node[side, anchor=south east] at (118,0.4) {\iflangja{メモリ}{memory}};
  \node[side, anchor=north east] at (118,-0.4) {\iflangja{ディスク}{disk}};
  % file descriptor table of process P
  \node[lab, anchor=south] at (9,42.2) {\iflangja{P の fd}{fds of P}};
  \foreach \i/\y in {0/40,1/36,2/32,3/28} {
    \draw[fill=white] (4,\y-2) rectangle (14,\y+2);
    \node[font=\scriptsize\ttfamily] at (9,\y) {\i};
  }
  \draw[fill=ln-accent!22] (4,26) rectangle (14,30);
  \node[font=\scriptsize\ttfamily] at (9,28) {3};
  % file
  \node[obj, minimum width=26mm] (file) at (36,28)
    {\iflangja{ファイル\\（オフセット，モード）}{file\\(offset, mode)}};
  \draw[->] (14,28) -- (file.west);
  % dentries
  \draw[densely dotted, ln-muted, fill=ln-box] (3,9) rectangle (70,19);
  \node[lab, anchor=north west, text=ln-muted] at (3.5,8.6) {\iflangja{dentry キャッシュ（メモリ上のみ）}{dentry cache (in memory only)}};
  \node[den] (r) at (10,14) {/};
  \node[den] (u) at (24,14) {users};
  \node[den] (j) at (39,14) {joe};
  \node[den] (n) at (58,14) {notes.txt};
  \draw[->] (r) -- (u);  \draw[->] (u) -- (j);  \draw[->] (j) -- (n);
  \draw[->] (file.east) -| (n.north);
  % inode in memory
  \node[obj, minimum width=32mm, minimum height=15mm] (ino) at (97,20)
    {\iflangja{iノード\\所有者，権限，サイズ，\\ブロックの一覧}{inode\\owner, permissions, size,\\list of blocks}};
  \draw[->] (n.east) -- (ino.west |- n.east);
  % on-disk layout
  \def\seg#1#2#3#4{\draw[fill=#3] (#1,-15) rectangle (#2,-8);
    \node[lab] at ({(#1+#2)/2},-11.5) {#4};}
  \seg{4}{30}{ln-accent!22}{\iflangja{スーパーブロック}{superblock}}
  \seg{30}{58}{ln-box}{\iflangja{iノードブロック}{inode blocks}}
  \seg{58}{98}{white}{\iflangja{データブロック}{data blocks}}
  \seg{98}{118}{black!8}{\iflangja{空きブロック}{free blocks}}
  \draw[<->, densely dashed] (85,12.5) -- (48,-8);
  \node[lab, anchor=east] at (64,3.2) {\iflangja{ディスク上の表現}{on-disk copy}};
  \foreach \x in {66,78,90} \draw[->] (97,12.5) -- (\x,-8);
  \draw[->, ln-accent] (17,-8) .. controls (22,-4) and (36,-4) .. (44,-8);
  \node[lab, text=ln-accent, anchor=south] at (30,-3.7) {\iflangja{配置の地図}{map of the layout}};
\end{tikzpicture}

  \caption{VFS abstractions.  File descriptor~3 of process~P refers to an
    open file, whose dentry leads to the inode; the dentries of the path
    exist only in memory.  On disk, the superblock maps the blocks that hold
    inodes, data and free space.}
  \label{fig:l11-vfs-objects}
\end{figure}

\begin{definition}[File descriptor]
  A \term{file descriptor}[ファイルディスクリプタ] is the operating system's
  representation of an open file for a process: a small integer that the
  process obtains from \texttt{open} and passes to the operations on the
  file, such as \texttt{read}, \texttt{write}, \texttt{sendfile},
  locking and \texttt{close}.
\end{definition}

The properties of the file itself, and the locations of its data blocks,
are recorded in its inode.

\begin{definition}[Inode]
  An \term{inode}[iノード] is the persistent representation of a file: an
  index that lists all the data blocks of the file, together with metadata
  such as the device, the permissions and the size of the file.
\end{definition}

A directory is a file, too; its data consist of information about the files
it contains and their inodes.  File names therefore appear in directories,
not in inodes.  The components of a path are represented by further objects.

\begin{definition}[Directory entry]
  A \term{directory entry}[ディレクトリエントリ]%
  \index{dentry|see{directory entry}} (dentry) is the VFS object that
  corresponds to a single component of a path.
\end{definition}

Looking up the path \texttt{/users/joe} creates a dentry for each of
\texttt{/}, \texttt{users} and \texttt{joe}.  The VFS keeps dentries in the
\emph{dentry cache}, so that a later lookup of a path with the same prefix
does not have to read the corresponding directories again.  Dentries are
soft state: they exist only in memory and have no on-disk representation.

\begin{definition}[Superblock]
  A \term{superblock}[スーパーブロック] contains file-system-specific
  information about the layout of a file system.
\end{definition}

\begin{example}
  A process opens \texttt{/users/joe/notes.txt}.  Starting from the dentry
  of \texttt{/}, the VFS reads each directory on the path to find the next
  component and its inode, creating the dentries of \texttt{users},
  \texttt{joe} and \texttt{notes.txt} on the way.  If the process then opens
  \texttt{/users/joe/todo.txt}, the dentries of \texttt{/}, \texttt{users}
  and \texttt{joe} are found in the dentry cache, and only the last
  component must be looked up in the directory \texttt{joe}.
\end{example}

\subsection{VFS on disk}

On disk, the abstractions are placed as follows (\cref{fig:l11-vfs-objects}).
The data of a file are kept in data blocks.  The inode, which tracks the
blocks of the file, also resides on disk, in some block.  The superblock
provides the overall map of the disk blocks: which blocks hold inodes, which
hold data, and which are free.

\section{The second extended file system}
\label{sec:l11-ext2}

Linux has long used a family of file systems whose on-disk structures
correspond closely to the abstractions of the VFS.\source{P3L5, ext2;
OSTEP ch.~40--41.}  The second extended file system, \term{ext2}[ext2], was
for many years the default file system of Linux and is the basis of its
successors ext3 and ext4.  An ext2 partition begins with \SI{1024}{\byte}
reserved for booting---all of block~0 when blocks are
\SI{1}{\kibi\byte}---and the rest is divided into groups of blocks
(\cref{fig:l11-ext2}).

\begin{definition}[Block group]
  A \term{block group}[ブロックグループ] is one of the equal-sized groups of
  consecutive blocks into which ext2 divides a disk partition.  Each block
  group holds its own inodes and data blocks together with the metadata that
  describe them.
\end{definition}

Each block group contains
\begin{description}
  \item[a copy of the superblock] recording file-system-wide information:
    the numbers of inodes and of disk blocks, the numbers of free inodes and
    blocks, and the size of a block;
  \item[a copy of the group descriptor table] with one descriptor per group,
    recording the locations of the bitmaps and the inode table of that group
    and its numbers of free blocks, free inodes and directories;
  \item[a block bitmap and an inode bitmap] tracking which blocks and which
    inodes of the group are free;
  \item[the inode table] holding a consecutive range of the inodes of the
    file system, which are numbered from 1 to the maximum number, one per
    file;
  \item[the data blocks] holding the file data.
\end{description}
The superblock and the group descriptor table describe the file system as a
whole; the copies in block group~0 are the ones normally used, and the
copies in the other groups serve as backups.  Keeping the inodes and the
data blocks of a file in the same group places related blocks close together
on the disk and reduces the distance the disk head must travel between
them.\margin{A group is not a region of the disk
geometry: it holds as many blocks as one block of bitmap can track.}

\begin{figure}[htbp]
  \centering
  % Layout of an ext2 partition: boot block and block groups; contents of one
% block group.
% Usage: % Layout of an ext2 partition: boot block and block groups; contents of one
% block group.
% Usage: % Layout of an ext2 partition: boot block and block groups; contents of one
% block group.
% Usage: \input{figures/l11-ext2}   (width ~116 mm)
\begin{tikzpicture}[x=1mm, y=1mm, font=\scriptsize\sffamily,
  lab/.style={font=\scriptsize\sffamily, align=center},
  note/.style={font=\scriptsize\sffamily, text=ln-muted, align=center,
               anchor=north, text width=17mm, inner sep=0.5pt}]
  % top strip
  \def\tseg#1#2#3#4{\draw[fill=#3] (#1,18) rectangle (#2,25);
    \node[lab] at ({(#1+#2)/2},21.5) {#4};}
  \tseg{2}{22}{black!8}{\iflangja{ブート}{boot}}
  \tseg{22}{50}{white}{\iflangja{ブロックグループ 0}{block group 0}}
  \tseg{50}{78}{ln-box}{\iflangja{ブロックグループ 1}{block group 1}}
  \tseg{78}{88}{white}{\ldots}
  \tseg{88}{116}{white}{\iflangja{ブロックグループ $N$}{block group $N$}}
  \node[lab, anchor=south] at (12,25.3) {\iflangja{ブロック 0}{block 0}};
  % zoom lines
  \draw[densely dashed, ln-muted] (50,18) -- (2,5);
  \draw[densely dashed, ln-muted] (78,18) -- (116,5);
  % block group detail
  \def\bseg#1#2#3#4{\draw[fill=#3] (#1,-4) rectangle (#2,5);
    \node[lab] at ({(#1+#2)/2},0.5) {#4};}
  \bseg{2}{21}{ln-accent!22}{\iflangja{スーパー\\ブロック}{superblock}}
  \bseg{21}{40}{ln-accent!10}{\iflangja{グループ\\記述子}{group\\descriptor}}
  \bseg{40}{58}{ln-box}{\iflangja{ブロック\\ビットマップ}{block\\bitmap}}
  \bseg{58}{76}{ln-box}{\iflangja{iノード\\ビットマップ}{inode\\bitmap}}
  \bseg{76}{94}{white}{\iflangja{iノード}{inodes}}
  \bseg{94}{116}{white}{\iflangja{データ\\ブロック}{data\\blocks}}
  \node[note] at (11.5,-5) {\iflangja{iノード数，ブロック数，空きブロックの開始位置}{no.\ of inodes and blocks, start of free blocks}};
  \node[note] at (30.5,-5) {\iflangja{ビットマップの位置，空き iノード数，ディレクトリ数}{bitmaps, free inodes, directories}};
  \node[note] at (49,-5) {\iflangja{空きブロックを記録}{tracks free blocks}};
  \node[note] at (67,-5) {\iflangja{空き iノードを記録}{tracks free inodes}};
  \node[note] at (85,-5) {\iflangja{1 から最大数まで，ファイルごとに 1 つ}{1 to max.\ number, one per file}};
  \node[note] at (105,-5) {\iflangja{ファイルのデータ}{file data}};
\end{tikzpicture}
   (width ~116 mm)
\begin{tikzpicture}[x=1mm, y=1mm, font=\scriptsize\sffamily,
  lab/.style={font=\scriptsize\sffamily, align=center},
  note/.style={font=\scriptsize\sffamily, text=ln-muted, align=center,
               anchor=north, text width=17mm, inner sep=0.5pt}]
  % top strip
  \def\tseg#1#2#3#4{\draw[fill=#3] (#1,18) rectangle (#2,25);
    \node[lab] at ({(#1+#2)/2},21.5) {#4};}
  \tseg{2}{22}{black!8}{\iflangja{ブート}{boot}}
  \tseg{22}{50}{white}{\iflangja{ブロックグループ 0}{block group 0}}
  \tseg{50}{78}{ln-box}{\iflangja{ブロックグループ 1}{block group 1}}
  \tseg{78}{88}{white}{\ldots}
  \tseg{88}{116}{white}{\iflangja{ブロックグループ $N$}{block group $N$}}
  \node[lab, anchor=south] at (12,25.3) {\iflangja{ブロック 0}{block 0}};
  % zoom lines
  \draw[densely dashed, ln-muted] (50,18) -- (2,5);
  \draw[densely dashed, ln-muted] (78,18) -- (116,5);
  % block group detail
  \def\bseg#1#2#3#4{\draw[fill=#3] (#1,-4) rectangle (#2,5);
    \node[lab] at ({(#1+#2)/2},0.5) {#4};}
  \bseg{2}{21}{ln-accent!22}{\iflangja{スーパー\\ブロック}{superblock}}
  \bseg{21}{40}{ln-accent!10}{\iflangja{グループ\\記述子}{group\\descriptor}}
  \bseg{40}{58}{ln-box}{\iflangja{ブロック\\ビットマップ}{block\\bitmap}}
  \bseg{58}{76}{ln-box}{\iflangja{iノード\\ビットマップ}{inode\\bitmap}}
  \bseg{76}{94}{white}{\iflangja{iノード}{inodes}}
  \bseg{94}{116}{white}{\iflangja{データ\\ブロック}{data\\blocks}}
  \node[note] at (11.5,-5) {\iflangja{iノード数，ブロック数，空きブロックの開始位置}{no.\ of inodes and blocks, start of free blocks}};
  \node[note] at (30.5,-5) {\iflangja{ビットマップの位置，空き iノード数，ディレクトリ数}{bitmaps, free inodes, directories}};
  \node[note] at (49,-5) {\iflangja{空きブロックを記録}{tracks free blocks}};
  \node[note] at (67,-5) {\iflangja{空き iノードを記録}{tracks free inodes}};
  \node[note] at (85,-5) {\iflangja{1 から最大数まで，ファイルごとに 1 つ}{1 to max.\ number, one per file}};
  \node[note] at (105,-5) {\iflangja{ファイルのデータ}{file data}};
\end{tikzpicture}
   (width ~116 mm)
\begin{tikzpicture}[x=1mm, y=1mm, font=\scriptsize\sffamily,
  lab/.style={font=\scriptsize\sffamily, align=center},
  note/.style={font=\scriptsize\sffamily, text=ln-muted, align=center,
               anchor=north, text width=17mm, inner sep=0.5pt}]
  % top strip
  \def\tseg#1#2#3#4{\draw[fill=#3] (#1,18) rectangle (#2,25);
    \node[lab] at ({(#1+#2)/2},21.5) {#4};}
  \tseg{2}{22}{black!8}{\iflangja{ブート}{boot}}
  \tseg{22}{50}{white}{\iflangja{ブロックグループ 0}{block group 0}}
  \tseg{50}{78}{ln-box}{\iflangja{ブロックグループ 1}{block group 1}}
  \tseg{78}{88}{white}{\ldots}
  \tseg{88}{116}{white}{\iflangja{ブロックグループ $N$}{block group $N$}}
  \node[lab, anchor=south] at (12,25.3) {\iflangja{ブロック 0}{block 0}};
  % zoom lines
  \draw[densely dashed, ln-muted] (50,18) -- (2,5);
  \draw[densely dashed, ln-muted] (78,18) -- (116,5);
  % block group detail
  \def\bseg#1#2#3#4{\draw[fill=#3] (#1,-4) rectangle (#2,5);
    \node[lab] at ({(#1+#2)/2},0.5) {#4};}
  \bseg{2}{21}{ln-accent!22}{\iflangja{スーパー\\ブロック}{superblock}}
  \bseg{21}{40}{ln-accent!10}{\iflangja{グループ\\記述子}{group\\descriptor}}
  \bseg{40}{58}{ln-box}{\iflangja{ブロック\\ビットマップ}{block\\bitmap}}
  \bseg{58}{76}{ln-box}{\iflangja{iノード\\ビットマップ}{inode\\bitmap}}
  \bseg{76}{94}{white}{\iflangja{iノード}{inodes}}
  \bseg{94}{116}{white}{\iflangja{データ\\ブロック}{data\\blocks}}
  \node[note] at (11.5,-5) {\iflangja{iノード数，ブロック数，空きブロックの開始位置}{no.\ of inodes and blocks, start of free blocks}};
  \node[note] at (30.5,-5) {\iflangja{ビットマップの位置，空き iノード数，ディレクトリ数}{bitmaps, free inodes, directories}};
  \node[note] at (49,-5) {\iflangja{空きブロックを記録}{tracks free blocks}};
  \node[note] at (67,-5) {\iflangja{空き iノードを記録}{tracks free inodes}};
  \node[note] at (85,-5) {\iflangja{1 から最大数まで，ファイルごとに 1 つ}{1 to max.\ number, one per file}};
  \node[note] at (105,-5) {\iflangja{ファイルのデータ}{file data}};
\end{tikzpicture}

  \caption{Layout of an ext2 partition.  After the space reserved for
    booting, the partition is divided into block groups, each with a copy of
    the superblock and of the group descriptor table, its own bitmaps, and
    its inodes and data blocks.}
  \label{fig:l11-ext2}
\end{figure}

\needspace{6\baselineskip}
\section{Inodes}
\label{sec:l11-inodes}

A file is identified by its inode, and the inode is the index of all disk
blocks of the file: it lists the blocks, in order, together with the
metadata.\source{P3L5, inodes, inodes with indirect pointers; OSTEP ch.~40;
Kerrisk ch.~14.}  In its simplest form the inode contains the block numbers
of all blocks of the file.  Sequential and random access are then equally
easy: with block size $B$, the byte at offset $o$ lies in the block whose
number is entry $\lfloor o/B \rfloor$ of the list.  The size of a file,
however, is limited by the number of entries that fit into the inode.

\begin{example}
  An inode of 512~bytes that uses 64~bytes for metadata has room
  for $448/4 = 112$ block numbers of 4~bytes each.  With
  \SI{4}{\kibi\byte} blocks, no file can be larger than
  $112 \times \SI{4}{\kibi\byte} = \SI{448}{\kibi\byte}$.
\end{example}

\subsection{Indirect pointers}

To support large files with small inodes, an inode holds pointers of several
kinds (\cref{fig:l11-inode}).
\begin{description}
  \item[Direct pointers] A \term{direct pointer}[直接ポインタ] points to a
    data block.
  \item[Indirect pointers] An \term{indirect pointer}[間接ポインタ] points
    to a block that contains pointers to data blocks.
  \item[Double indirect pointers] A
    \term{double indirect pointer}[二重間接ポインタ] points to a block of
    pointers to blocks of pointers to data blocks.
\end{description}
A \emph{triple indirect pointer} adds a further level.  The inodes of ext2
and of many UNIX file systems contain 12 direct pointers and one pointer of
each indirect kind.

\begin{figure}[htbp]
  \centering
  % An inode with direct, single, double and triple indirect pointers.
% Usage: % An inode with direct, single, double and triple indirect pointers.
% Usage: % An inode with direct, single, double and triple indirect pointers.
% Usage: \input{figures/l11-inode}   (width ~116 mm)
\begin{tikzpicture}[x=1mm, y=1mm, font=\scriptsize\sffamily,
  >={Stealth[length=1.4mm]},
  data/.style={draw, fill=white, minimum width=10mm, minimum height=4.5mm,
               inner sep=0.5pt, font=\scriptsize\sffamily},
  lab/.style={font=\scriptsize\sffamily, align=center}]
  % \ptrblk{name}{x}{y}: block of pointers
  \def\ptrblk#1#2#3{%
    \draw[fill=ln-box] (#2-5,#3-4) rectangle (#2+5,#3+4);
    \foreach \k in {-2,0,2} \draw[ln-rule] (#2-5,#3+\k) -- (#2+5,#3+\k);
    \coordinate (#1w) at (#2-5,#3);
    \coordinate (#1e) at (#2+5,#3+3);
    \coordinate (#1s) at (#2+5,#3-3);}
  % inode
  \node[font=\footnotesize\sffamily\bfseries, anchor=south] at (19,48.3) {\iflangja{iノード}{inode}};
  \draw[fill=ln-accent!10] (4,36) rectangle (34,48);
  \node[lab] at (19,42) {\iflangja{メタデータ\\（モード，所有者，\\タイムスタンプ，サイズ）}{metadata\\(mode, owners,\\timestamps, size, \ldots)}};
  \def\irow#1#2{\draw[fill=white] (4,#1-2.5) rectangle (34,#1+2.5);
    \node[lab] at (19,#1) {#2};}
  \irow{33.5}{\iflangja{直接ポインタ 0}{direct 0}}
  \irow{28.5}{\ldots}
  \irow{23.5}{\iflangja{直接ポインタ 11}{direct 11}}
  \irow{18.5}{\iflangja{間接ポインタ}{single indirect}}
  \irow{13.5}{\iflangja{二重間接ポインタ}{double indirect}}
  \irow{8.5}{\iflangja{三重間接ポインタ}{triple indirect}}
  % direct
  \node[data] (a) at (111,33.5) {\iflangja{データ}{data}};
  \node[data] (b) at (111,23.5) {\iflangja{データ}{data}};
  \draw[->] (31,33.5) -- (a.west);
  \draw[->] (31,23.5) -- (b.west);
  % single indirect
  \ptrblk{p}{53}{15}
  \node[data] (c) at (111,15) {\iflangja{データ}{data}};
  \draw[->] (31,18.5) -- (pw);
  \draw[->] (pe) -- (c.west);
  % double indirect
  \ptrblk{q}{53}{0}
  \ptrblk{r}{77}{0}
  \node[data] (d) at (111,0) {\iflangja{データ}{data}};
  \draw[->] (31,13.5) -- (qw);
  \draw[->] (qe) -- (rw);
  \draw[->] (re) -- (d.west);
  % triple indirect
  \ptrblk{s}{53}{-15}
  \ptrblk{t}{71}{-15}
  \ptrblk{u}{89}{-15}
  \node[data] (e) at (111,-15) {\iflangja{データ}{data}};
  \draw[->] (31,8.5) -- (sw);
  \draw[->] (se) -- (tw);
  \draw[->] (te) -- (uw);
  \draw[->] (ue) -- (e.west);
  % legend
  \draw[fill=ln-box] (4,-27) rectangle (10,-23);
  \node[lab, anchor=west] at (11,-25) {\iflangja{ポインタのブロック}{block of pointers}};
  \draw[fill=white] (44,-27) rectangle (50,-23);
  \node[lab, anchor=west] at (51,-25) {\iflangja{データブロック}{data block}};
\end{tikzpicture}
   (width ~116 mm)
\begin{tikzpicture}[x=1mm, y=1mm, font=\scriptsize\sffamily,
  >={Stealth[length=1.4mm]},
  data/.style={draw, fill=white, minimum width=10mm, minimum height=4.5mm,
               inner sep=0.5pt, font=\scriptsize\sffamily},
  lab/.style={font=\scriptsize\sffamily, align=center}]
  % \ptrblk{name}{x}{y}: block of pointers
  \def\ptrblk#1#2#3{%
    \draw[fill=ln-box] (#2-5,#3-4) rectangle (#2+5,#3+4);
    \foreach \k in {-2,0,2} \draw[ln-rule] (#2-5,#3+\k) -- (#2+5,#3+\k);
    \coordinate (#1w) at (#2-5,#3);
    \coordinate (#1e) at (#2+5,#3+3);
    \coordinate (#1s) at (#2+5,#3-3);}
  % inode
  \node[font=\footnotesize\sffamily\bfseries, anchor=south] at (19,48.3) {\iflangja{iノード}{inode}};
  \draw[fill=ln-accent!10] (4,36) rectangle (34,48);
  \node[lab] at (19,42) {\iflangja{メタデータ\\（モード，所有者，\\タイムスタンプ，サイズ）}{metadata\\(mode, owners,\\timestamps, size, \ldots)}};
  \def\irow#1#2{\draw[fill=white] (4,#1-2.5) rectangle (34,#1+2.5);
    \node[lab] at (19,#1) {#2};}
  \irow{33.5}{\iflangja{直接ポインタ 0}{direct 0}}
  \irow{28.5}{\ldots}
  \irow{23.5}{\iflangja{直接ポインタ 11}{direct 11}}
  \irow{18.5}{\iflangja{間接ポインタ}{single indirect}}
  \irow{13.5}{\iflangja{二重間接ポインタ}{double indirect}}
  \irow{8.5}{\iflangja{三重間接ポインタ}{triple indirect}}
  % direct
  \node[data] (a) at (111,33.5) {\iflangja{データ}{data}};
  \node[data] (b) at (111,23.5) {\iflangja{データ}{data}};
  \draw[->] (31,33.5) -- (a.west);
  \draw[->] (31,23.5) -- (b.west);
  % single indirect
  \ptrblk{p}{53}{15}
  \node[data] (c) at (111,15) {\iflangja{データ}{data}};
  \draw[->] (31,18.5) -- (pw);
  \draw[->] (pe) -- (c.west);
  % double indirect
  \ptrblk{q}{53}{0}
  \ptrblk{r}{77}{0}
  \node[data] (d) at (111,0) {\iflangja{データ}{data}};
  \draw[->] (31,13.5) -- (qw);
  \draw[->] (qe) -- (rw);
  \draw[->] (re) -- (d.west);
  % triple indirect
  \ptrblk{s}{53}{-15}
  \ptrblk{t}{71}{-15}
  \ptrblk{u}{89}{-15}
  \node[data] (e) at (111,-15) {\iflangja{データ}{data}};
  \draw[->] (31,8.5) -- (sw);
  \draw[->] (se) -- (tw);
  \draw[->] (te) -- (uw);
  \draw[->] (ue) -- (e.west);
  % legend
  \draw[fill=ln-box] (4,-27) rectangle (10,-23);
  \node[lab, anchor=west] at (11,-25) {\iflangja{ポインタのブロック}{block of pointers}};
  \draw[fill=white] (44,-27) rectangle (50,-23);
  \node[lab, anchor=west] at (51,-25) {\iflangja{データブロック}{data block}};
\end{tikzpicture}
   (width ~116 mm)
\begin{tikzpicture}[x=1mm, y=1mm, font=\scriptsize\sffamily,
  >={Stealth[length=1.4mm]},
  data/.style={draw, fill=white, minimum width=10mm, minimum height=4.5mm,
               inner sep=0.5pt, font=\scriptsize\sffamily},
  lab/.style={font=\scriptsize\sffamily, align=center}]
  % \ptrblk{name}{x}{y}: block of pointers
  \def\ptrblk#1#2#3{%
    \draw[fill=ln-box] (#2-5,#3-4) rectangle (#2+5,#3+4);
    \foreach \k in {-2,0,2} \draw[ln-rule] (#2-5,#3+\k) -- (#2+5,#3+\k);
    \coordinate (#1w) at (#2-5,#3);
    \coordinate (#1e) at (#2+5,#3+3);
    \coordinate (#1s) at (#2+5,#3-3);}
  % inode
  \node[font=\footnotesize\sffamily\bfseries, anchor=south] at (19,48.3) {\iflangja{iノード}{inode}};
  \draw[fill=ln-accent!10] (4,36) rectangle (34,48);
  \node[lab] at (19,42) {\iflangja{メタデータ\\（モード，所有者，\\タイムスタンプ，サイズ）}{metadata\\(mode, owners,\\timestamps, size, \ldots)}};
  \def\irow#1#2{\draw[fill=white] (4,#1-2.5) rectangle (34,#1+2.5);
    \node[lab] at (19,#1) {#2};}
  \irow{33.5}{\iflangja{直接ポインタ 0}{direct 0}}
  \irow{28.5}{\ldots}
  \irow{23.5}{\iflangja{直接ポインタ 11}{direct 11}}
  \irow{18.5}{\iflangja{間接ポインタ}{single indirect}}
  \irow{13.5}{\iflangja{二重間接ポインタ}{double indirect}}
  \irow{8.5}{\iflangja{三重間接ポインタ}{triple indirect}}
  % direct
  \node[data] (a) at (111,33.5) {\iflangja{データ}{data}};
  \node[data] (b) at (111,23.5) {\iflangja{データ}{data}};
  \draw[->] (31,33.5) -- (a.west);
  \draw[->] (31,23.5) -- (b.west);
  % single indirect
  \ptrblk{p}{53}{15}
  \node[data] (c) at (111,15) {\iflangja{データ}{data}};
  \draw[->] (31,18.5) -- (pw);
  \draw[->] (pe) -- (c.west);
  % double indirect
  \ptrblk{q}{53}{0}
  \ptrblk{r}{77}{0}
  \node[data] (d) at (111,0) {\iflangja{データ}{data}};
  \draw[->] (31,13.5) -- (qw);
  \draw[->] (qe) -- (rw);
  \draw[->] (re) -- (d.west);
  % triple indirect
  \ptrblk{s}{53}{-15}
  \ptrblk{t}{71}{-15}
  \ptrblk{u}{89}{-15}
  \node[data] (e) at (111,-15) {\iflangja{データ}{data}};
  \draw[->] (31,8.5) -- (sw);
  \draw[->] (se) -- (tw);
  \draw[->] (te) -- (uw);
  \draw[->] (ue) -- (e.west);
  % legend
  \draw[fill=ln-box] (4,-27) rectangle (10,-23);
  \node[lab, anchor=west] at (11,-25) {\iflangja{ポインタのブロック}{block of pointers}};
  \draw[fill=white] (44,-27) rectangle (50,-23);
  \node[lab, anchor=west] at (51,-25) {\iflangja{データブロック}{data block}};
\end{tikzpicture}

  \caption{An inode with 12 direct pointers and single, double and triple
    indirect pointers.  Each level of indirection multiplies the number of
    reachable data blocks by the number of pointers per block.}
  \label{fig:l11-inode}
\end{figure}

With block size $B$ and pointer size $p$, a block holds $P = B/p$ pointers.
An inode with $n_d$ direct pointers and one pointer of each of the three
indirect kinds can address files of up to
\begin{equation}
  S_{\max} \;=\; \bigl(n_d + P + P^2 + P^3\bigr) \times B .
  \label{eq:l11-max-size}
\end{equation}

\begin{example}
  With \SI{4}{\kibi\byte} blocks and 8~bytes per pointer, a block
  holds $P = 4096/8 = 512$ pointers.  The 12 direct pointers address
  $12 \times \SI{4}{\kibi\byte} = \SI{48}{\kibi\byte}$; the indirect pointer
  $512 \times \SI{4}{\kibi\byte} = \SI{2}{\mebi\byte}$; the double indirect
  pointer $512^2 \times \SI{4}{\kibi\byte} = \SI{1}{\gibi\byte}$; and the
  triple indirect pointer $512^3 \times \SI{4}{\kibi\byte} =
  \SI{512}{\gibi\byte}$.  The maximum file size is about
  \SI{513}{\gibi\byte}.
\end{example}

The benefit of indirect pointers is that a small inode supports very large
files.  The cost is slower file access, because every level of indirection
adds a disk access to reach a data block.  Reading a block through a direct
pointer requires two disk accesses, one for the inode and one for the data
block; reading a block through a double indirect pointer requires up to
four.

\begin{example}
  With the parameters of the previous example, a process reads the byte at
  offset \SI{600}{\mebi\byte} of a file.  It lies in block
  $600 \times 1024 / 4 = \num{153600}$ of the file.  Blocks 0--11 are reached
  directly, blocks 12--523 through the indirect pointer, and blocks
  \num{524}--\num{262667} through the double indirect pointer, so the read
  requires the inode, two blocks of pointers and the data block: four disk
  accesses if none of them is cached.
\end{example}

\begin{keypoint}
  The VFS operates on files through file descriptors, dentries, inodes and
  superblocks.  On disk, a file is its inode plus the data blocks the inode
  indexes; direct and indirect pointers let a fixed-size inode index very
  large files at the price of extra disk accesses per level of indirection.
\end{keypoint}

\section{Disk access optimizations}
\label{sec:l11-optim}

Accessing a file costs disk accesses, and on a hard disk every access that
is not sequential costs time to move the head.\source{P3L5, disk access
optimizations; OSTEP ch.~37, 40, 42; Kerrisk ch.~13--14.}  File systems use
four techniques to reduce the overhead of file access
(\cref{tab:l11-optim}).

\begin{table}[htbp]
  \centering
  \caption{Techniques that reduce file access overhead.}
  \label{tab:l11-optim}
  \small
  \begin{tabular}{@{}l>{\raggedright\arraybackslash}p{0.62\linewidth}@{}}
    \toprule
    Technique & Effect \\
    \midrule
    Caching and buffering & fewer disk accesses \\
    I/O scheduling & less head movement; more sequential, less random
      access \\
    Prefetching & more cache hits, by exploiting locality \\
    Journaling & less random access for writes, without keeping updates
      only in memory \\
    \bottomrule
  \end{tabular}
\end{table}

\subsection{Caching and buffering}

\begin{definition}[Buffer cache]
  A \term{buffer cache}[バッファキャッシュ] is a cache of file data in main
  memory.  Reads are served from the cache when possible and writes update
  the cache; modified data are written to disk periodically, or when a
  process requests it with \texttt{fsync}.
\end{definition}

A process that must make an update durable follows the write with
\texttt{fsync}:
\needspace{5\baselineskip}
\begin{lstlisting}[language=C, numbers=none]
write(fd, rec, sizeof *rec);  /* updates the cache */
fsync(fd);                    /* flushes it to disk */
\end{lstlisting}
\caution{A successful \texttt{write} does not mean the data are on disk; a
crash loses data still in the cache.}
\texttt{fsync} returns only after the data of the file have been written to
the device.

\subsection{I/O scheduling}

Ordering the pending disk requests so as to reduce the movement of the disk
head is called \term{I/O scheduling}[I/Oスケジューリング]: requests are
served in an order that maximizes sequential access and minimizes random
access.

\begin{example}
  The head is positioned at block 10, and writes to blocks 90, 12, 91, 13
  and 50 arrive in this order.  Assuming that block numbers reflect
  positions on the disk, serving the requests in arrival order moves the
  head across $80+78+79+78+37 = 352$ blocks.  Serving them in the order 12,
  13, 50, 90, 91 moves it across only $2+1+37+40+1 = 81$ blocks, and the
  pairs 12, 13 and 90, 91 become sequential writes.
\end{example}

\needspace{5\baselineskip}
\subsection{Prefetching}

Reading more blocks than a request asks for is called
\term{prefetching}[プリフェッチ] or read-ahead: when a block is read, the
blocks likely to be needed next are read with it.  Prefetching exploits locality---a file that
is being read sequentially will soon need the following blocks---and
increases the hit rate of the buffer cache.  If a process reads block 40 of
a file and the file system also reads blocks 41--43 in the same disk
operation, the next three reads are served from memory without further
head movement.

\subsection{Journaling}

Writing every update to its proper location immediately causes random
access, whereas keeping updates only in the buffer cache risks losing them.
The technique of \term{journaling}[ジャーナリング], or logging, avoids both
problems.  The file system describes each update---which block, at which
offset, with which value---in a log, the \emph{journal}, which it writes sequentially to a
dedicated area of the disk, and it periodically applies the logged updates
to their proper disk locations (\cref{fig:l11-journal}).  Because the
journal is on disk, an update recorded in it survives a crash and can be
applied afterwards.  The ext3 and ext4 file systems add such a journal to
the design of ext2; by default they log only metadata blocks and write file
data directly to their final locations, and logging the data as well is an
option.

\begin{figure}[htbp]
  \centering
  % Journaling: updates are appended sequentially to the journal and later
% applied to their proper locations on disk.
% Usage: % Journaling: updates are appended sequentially to the journal and later
% applied to their proper locations on disk.
% Usage: % Journaling: updates are appended sequentially to the journal and later
% applied to their proper locations on disk.
% Usage: \input{figures/l11-journal}   (width ~116 mm)
\begin{tikzpicture}[x=1mm, y=1mm, font=\scriptsize\sffamily,
  >={Stealth[length=1.5mm]},
  lab/.style={font=\scriptsize\sffamily, align=center},
  num/.style={circle, draw=black, fill=white, text=black, line width=0.4pt,
              inner sep=0.4pt, font=\scriptsize\sffamily\bfseries}]
  % updates
  \node[draw, fill=white, align=center, inner sep=1.5pt] (up) at (39,22)
    {\iflangja{更新：ブロック 870，12，301}{updates to blocks 870, 12, 301}};
  \draw[->, line width=1.2pt, ln-accent] (up.south) -- node[num, right=0.8mm] {1} (39,8.2);
  \node[lab, anchor=west] at (43,14) {\iflangja{順に追記（シーケンシャル書き込み）}{appended in order (sequential writes)}};
  % disk strip
  \node[lab, anchor=east] at (17,4) {\iflangja{ディスク}{disk}};
  \draw[fill=white] (18,0) rectangle (116,8);
  \foreach \x/\b in {18/870, 32/12, 46/301} {
    \draw[fill=ln-accent!10] (\x,0) rectangle (\x+14,8);
    \node[font=\scriptsize\ttfamily] at (\x+7,4) {\b};
  }
  \foreach \x/\b in {66/12, 84/301, 102/870} {
    \draw[fill=ln-box] (\x,0) rectangle (\x+10,8);
    \node[font=\scriptsize\ttfamily] at (\x+5,4) {\b};
  }
  \node[lab, anchor=south west, text=ln-muted] at (18,8.3) {\iflangja{ジャーナル}{journal}};
  \node[lab, anchor=south, text=ln-muted] at (89,8.3) {\iflangja{本来の位置}{proper locations}};
  % apply later
  \draw[->, densely dashed] (25,0) .. controls (30,-15) and (100,-15) .. (107,0);
  \draw[->, densely dashed] (39,0) .. controls (42,-7) and (68,-7) .. (71,0);
  \draw[->, densely dashed] (53,0) .. controls (58,-10) and (84,-10) .. (89,0);
  \node[num] at (66,-11.3) {2};
  \node[lab, anchor=west] at (69,-13.5) {\iflangja{後でまとめて反映}{applied later}};
\end{tikzpicture}
   (width ~116 mm)
\begin{tikzpicture}[x=1mm, y=1mm, font=\scriptsize\sffamily,
  >={Stealth[length=1.5mm]},
  lab/.style={font=\scriptsize\sffamily, align=center},
  num/.style={circle, draw=black, fill=white, text=black, line width=0.4pt,
              inner sep=0.4pt, font=\scriptsize\sffamily\bfseries}]
  % updates
  \node[draw, fill=white, align=center, inner sep=1.5pt] (up) at (39,22)
    {\iflangja{更新：ブロック 870，12，301}{updates to blocks 870, 12, 301}};
  \draw[->, line width=1.2pt, ln-accent] (up.south) -- node[num, right=0.8mm] {1} (39,8.2);
  \node[lab, anchor=west] at (43,14) {\iflangja{順に追記（シーケンシャル書き込み）}{appended in order (sequential writes)}};
  % disk strip
  \node[lab, anchor=east] at (17,4) {\iflangja{ディスク}{disk}};
  \draw[fill=white] (18,0) rectangle (116,8);
  \foreach \x/\b in {18/870, 32/12, 46/301} {
    \draw[fill=ln-accent!10] (\x,0) rectangle (\x+14,8);
    \node[font=\scriptsize\ttfamily] at (\x+7,4) {\b};
  }
  \foreach \x/\b in {66/12, 84/301, 102/870} {
    \draw[fill=ln-box] (\x,0) rectangle (\x+10,8);
    \node[font=\scriptsize\ttfamily] at (\x+5,4) {\b};
  }
  \node[lab, anchor=south west, text=ln-muted] at (18,8.3) {\iflangja{ジャーナル}{journal}};
  \node[lab, anchor=south, text=ln-muted] at (89,8.3) {\iflangja{本来の位置}{proper locations}};
  % apply later
  \draw[->, densely dashed] (25,0) .. controls (30,-15) and (100,-15) .. (107,0);
  \draw[->, densely dashed] (39,0) .. controls (42,-7) and (68,-7) .. (71,0);
  \draw[->, densely dashed] (53,0) .. controls (58,-10) and (84,-10) .. (89,0);
  \node[num] at (66,-11.3) {2};
  \node[lab, anchor=west] at (69,-13.5) {\iflangja{後でまとめて反映}{applied later}};
\end{tikzpicture}
   (width ~116 mm)
\begin{tikzpicture}[x=1mm, y=1mm, font=\scriptsize\sffamily,
  >={Stealth[length=1.5mm]},
  lab/.style={font=\scriptsize\sffamily, align=center},
  num/.style={circle, draw=black, fill=white, text=black, line width=0.4pt,
              inner sep=0.4pt, font=\scriptsize\sffamily\bfseries}]
  % updates
  \node[draw, fill=white, align=center, inner sep=1.5pt] (up) at (39,22)
    {\iflangja{更新：ブロック 870，12，301}{updates to blocks 870, 12, 301}};
  \draw[->, line width=1.2pt, ln-accent] (up.south) -- node[num, right=0.8mm] {1} (39,8.2);
  \node[lab, anchor=west] at (43,14) {\iflangja{順に追記（シーケンシャル書き込み）}{appended in order (sequential writes)}};
  % disk strip
  \node[lab, anchor=east] at (17,4) {\iflangja{ディスク}{disk}};
  \draw[fill=white] (18,0) rectangle (116,8);
  \foreach \x/\b in {18/870, 32/12, 46/301} {
    \draw[fill=ln-accent!10] (\x,0) rectangle (\x+14,8);
    \node[font=\scriptsize\ttfamily] at (\x+7,4) {\b};
  }
  \foreach \x/\b in {66/12, 84/301, 102/870} {
    \draw[fill=ln-box] (\x,0) rectangle (\x+10,8);
    \node[font=\scriptsize\ttfamily] at (\x+5,4) {\b};
  }
  \node[lab, anchor=south west, text=ln-muted] at (18,8.3) {\iflangja{ジャーナル}{journal}};
  \node[lab, anchor=south, text=ln-muted] at (89,8.3) {\iflangja{本来の位置}{proper locations}};
  % apply later
  \draw[->, densely dashed] (25,0) .. controls (30,-15) and (100,-15) .. (107,0);
  \draw[->, densely dashed] (39,0) .. controls (42,-7) and (68,-7) .. (71,0);
  \draw[->, densely dashed] (53,0) .. controls (58,-10) and (84,-10) .. (89,0);
  \node[num] at (66,-11.3) {2};
  \node[lab, anchor=west] at (69,-13.5) {\iflangja{後でまとめて反映}{applied later}};
\end{tikzpicture}

  \caption{Journaling.  Updates to scattered blocks are first appended in
    order to the journal and later applied to their proper locations.}
  \label{fig:l11-journal}
\end{figure}

\needspace{16\baselineskip}
\begin{keypoint}[Lesson summary]
  An I/O device presents command, data and status registers, reached
  through interconnects such as PCIe and operated by a device driver behind
  an interface standardized by the operating system.  Devices are block,
  character or network devices and are represented as device files.  The
  CPU addresses registers through memory-mapped I/O or I/O ports, learns of
  device events through interrupts or polling, and transfers data with
  programmed I/O or DMA.  A device access normally passes through the
  kernel, but operating system bypass can remove the kernel from the data
  path; the requesting thread either blocks or continues asynchronously.
  Files reach block devices through the file system, the generic block layer
  and the driver.  The virtual file system unifies file systems through file
  descriptors, dentries, inodes and superblocks; ext2 organizes a partition
  into block groups, and its inodes use direct and indirect pointers.
  Caching, I/O scheduling, prefetching and journaling reduce the cost of
  disk access.
\end{keypoint}

\end{document}
